\chapter{Zupełność modelowa}
\label{ch:mod-model-completeness}
\section{Motywacja i podstawowe pojęcia}
Eliminacja kwantyfikatorów pozwala badać wszystkie formuły
za pomocą atomów i spójników. W wielu teoriach wystarczy
jednak słabsza własność: aby zanurzenie jednego modelu
w drugi było elementarne, nie musimy eliminować wszystkich
kwantyfikatorów, lecz wystarczy odpowiednio kontrolować
formuły egzystencjalne.

W całym rozdziale pracujemy w logice pierwszego rzędu
z równością i w ZFC. Uniwersa struktur są niepuste.
Stosujemy konwencję stałych logicznych $\top,\bot$
z rozdziału \ref{ch:mod-qe}.
Przez zanurzenie rozumiemy izomorfizm na podstrukturę
obrazu (z zachowaniem i odbijaniem relacji), a nie
dowolny różnowartościowy homomorfizm. Pojęcia
podstruktury i elementarności pochodzą z definicji
\ref{def:model-substructure} oraz
\ref{def:model-elementary-substructure}.

\begin{definicja}[Zupełność modelowa]
\label{def:mod-model-complete}
Spełnialną teorię $T$ nazywamy modelowo zupełną, jeżeli
dla dowolnych modeli $\mathfrak M,\mathfrak N\models T$
każde zanurzenie $\mathfrak M\to\mathfrak N$ jest
elementarne. Równoważnie, jeżeli
$\mathfrak M\subseteq\mathfrak N$ są modelami $T$,
to $\mathfrak M\preccurlyeq\mathfrak N$.
\end{definicja}

Równoważność obu zapisów wynika z niezmienniczości
prawdziwości formuł względem izomorfizmów
(twierdzenie \ref{thm:model-iso-invariance}):
dowolne zanurzenie jest złożeniem izomorfizmu na obraz
i inkluzji obrazu w strukturę docelową.

Zupełność modelowa nie jest tym samym co zupełność teorii.
Zupełność dotyczy rozstrzygania zdań, natomiast zupełność
modelowa dotyczy zanurzeń oraz formuł z parametrami.
Nie należy też mylić obu własności z twierdzeniem
o pełności rachunku logicznego.

\begin{wniosek}
\label{cor:mod-mc-qe}
Każda teoria z eliminacją kwantyfikatorów jest
modelowo zupełna.
\end{wniosek}
\begin{proof}
Jest to pierwsza część wniosku
\ref{cor:mod-qe-consequences}: podstruktura zachowuje
formuły bez kwantyfikatorów, a eliminacja przenosi
tę zgodność na wszystkie formuły.
\end{proof}

\begin{przyklad}[Zupełność modelowa bez zupełności]
\label{ex:mod-mc-incomplete}
Teoria wszystkich ciał algebraicznie domkniętych
w języku pierścieni jest modelowo zupełna na mocy
twierdzenia \ref{thm:mod-qe-acf} i wniosku
\ref{cor:mod-mc-qe}. Nie jest jednak zupełna:
zdanie $1+1=0$ zachodzi w algebraicznym domknięciu
ciała dwuelementowego, lecz nie zachodzi w $\mathbb C$.
\end{przyklad}

\begin{przyklad}[Zupełność bez zupełności modelowej]
\label{ex:mod-complete-not-mc}
Teoria $\operatorname{Th}(\mathbb Z,<)$ jest zupełna,
ale nie jest modelowo zupełna. Odwzorowanie
$f:\mathbb Z\to\mathbb Z$, $f(n)=2n$, jest zanurzeniem
porządków liniowych. W dziedzinie formuła
$\exists y\,(x<y\land y<z)$ jest fałszywa przy $(0,1)$,
a w przeciwdziedzinie jest prawdziwa przy $(f(0),f(1))=(0,2)$,
ze świadkiem $1$. Zanurzenie nie jest więc elementarne.
\end{przyklad}

\section{Formuły egzystencjalne i domknięcie egzystencjalne}
\begin{definicja}[Formuły egzystencjalne i uniwersalne]
\label{def:mod-mc-existential-universal}
Formułę nazywamy egzystencjalną, jeśli ma postać
$\exists\bar y\,\theta(\bar x,\bar y)$, gdzie $\theta$
nie zawiera kwantyfikatorów, a uniwersalną, jeśli ma
postać $\forall\bar y\,\theta(\bar x,\bar y)$.
Lista $\bar y$ może być pusta.
W obu przypadkach zmienne związane są parami różne
i różne od wskazanych zmiennych wolnych $\bar x$.
Lista $\bar x$ może być pusta; obejmuje wszystkie
zmienne wolne rozważanej formuły, ale nie każda jej
zmienna musi wystąpić w macierzy.
\end{definicja}

W macierzy dopuszczamy negacje atomów, w przeciwieństwie
do formuł pozytywnych egzystencjalnych, omawianych
w definicji \ref{def:model-positive-existential}.
Zachowanie tych ostatnich przez homomorfizmy nie jest
tym samym co zachowanie wszystkich formuł egzystencjalnych
przez zanurzenia.

\begin{lemat}[Zachowanie przy inkluzjach]
\label{lem:mod-mc-preservation}
Niech $\mathfrak M\subseteq\mathfrak N$ i $\bar a\in M^n$.
Jeżeli $\varphi(\bar x)$ jest egzystencjalna, to
\[
\mathfrak M\models\varphi(\bar a)
\ \Longrightarrow\
\mathfrak N\models\varphi(\bar a).
\]
Jeżeli $\psi(\bar x)$ jest uniwersalna, to
\[
\mathfrak N\models\psi(\bar a)
\ \Longrightarrow\
\mathfrak M\models\psi(\bar a).
\]
\end{lemat}
\begin{proof}
Wartości termów i prawdziwość atomów przy argumentach
z $M$ są zgodne w obu strukturach na mocy definicji
podstruktury. Spójniki zachowują zgodność, więc
macierz bez kwantyfikatorów ma w obu strukturach
tę samą wartość przy takich argumentach.
Każdy świadek egzystencjalny w $M$ jest zatem świadkiem
w $N$. Dla formuły uniwersalnej ograniczamy wszystkie
wartości zmiennych związanych z $N$ do $M$.
\end{proof}

\begin{definicja}[Podstruktura egzystencjalnie domknięta]
\label{def:mod-mc-existential-closure}
Piszemy $\mathfrak M\preccurlyeq_{\exists}\mathfrak N$,
jeżeli $\mathfrak M\subseteq\mathfrak N$ i dla każdej
formuły egzystencjalnej $\varphi(\bar x)$ oraz każdej
krotki $\bar a$ z $M$ zachodzi
\[
\mathfrak N\models\varphi(\bar a)
\ \Longrightarrow\
\mathfrak M\models\varphi(\bar a).
\]
Mówimy wtedy, że $\mathfrak M$ jest egzystencjalnie
domknięta w $\mathfrak N$.
Jeżeli $\mathfrak M$ jest modelem teorii $T$ i
$\mathfrak M\preccurlyeq_{\exists}\mathfrak N$ dla każdego
rozszerzenia $\mathfrak N\supseteq\mathfrak M$
będącego modelem $T$, nazywamy
$\mathfrak M$ egzystencjalnie domkniętym modelem $T$.
\end{definicja}

Z lematu \ref{lem:mod-mc-preservation} wynika, że
$\preccurlyeq_{\exists}$ oznacza zgodność prawdziwości
wszystkich formuł egzystencjalnych przy parametrach
z mniejszej struktury. Jest też równoważna zgodności
wszystkich formuł uniwersalnych, ponieważ negacja
formuły uniwersalnej jest równoważna egzystencjalnej
z zanegowaną macierzą i odwrotnie.

\section{Diagramy i kryterium Robinsona}
W dowodach wykorzystamy diagramy struktur. Diagram
elementarny omawiano w związku z realizacją typów
w rozszerzeniach elementarnych (twierdzenie
\ref{thm:mod-type-elementary-extension}).

\begin{definicja}[Diagram bez kwantyfikatorów i diagram elementarny]
\label{def:mod-mc-diagrams}
Dla struktury $\mathfrak M$ rozszerzamy język o świeżą
stałą $c_a$ dla każdego $a\in M$, interpretowaną jako $a$.
Diagram bez kwantyfikatorów $\operatorname{Diag}(\mathfrak M)$
jest zbiorem wszystkich zdań bez kwantyfikatorów
prawdziwych w tak rozszerzonej strukturze.
Diagram elementarny $\operatorname{ED}(\mathfrak M)$
jest zbiorem wszystkich prawdziwych zdań rozszerzonego
języka, bez ograniczenia na kwantyfikatory.
\end{definicja}

\begin{lemat}[Własność diagramów]
\label{lem:mod-mc-diagram-embedding}
Jeśli $\mathfrak D$ jest modelem diagramu
$\operatorname{Diag}(\mathfrak M)$, odwzorowanie
$a\mapsto c_a^{\mathfrak D}$ jest zanurzeniem
$\mathfrak M$ w redukt $\mathfrak D$ do języka pierwotnego.
Jeśli $\mathfrak D$ jest modelem
$\operatorname{ED}(\mathfrak M)$, to zanurzenie
jest elementarne. Odwrotnie każde zanurzenie
(odpowiednio elementarne) daje model właściwego diagramu
przez interpretowanie nowych stałych obrazami elementów.
\end{lemat}
\begin{proof}
Dla różnych $a,b$ diagram zawiera $c_a\ne c_b$,
więc odwzorowanie jest różnowartościowe. Dla każdej
stałej języka pierwotnego zawiera równość z nazwą
jej wartości w $\mathfrak M$. Dla symbolu funkcyjnego
$f$ i argumentów $\bar a$ zawiera równość
$f(c_{\bar a})=c_{f^{\mathfrak M}(\bar a)}$.
Dla symbolu relacyjnego zawiera atom lub jego negację
zależnie od prawdziwości w $\mathfrak M$. Są to dokładnie
warunki izomorfizmu na podstrukturę obrazu.
W przypadku symbolu relacyjnego o arności zero,
jeżeli język go dopuszcza, diagram ustala bezpośrednio
jego wartość logiczną. Zgodność wszystkich atomów
daje zgodność każdej formuły bez kwantyfikatorów,
przez indukcję po spójnikach.

W przypadku diagramu elementarnego dla każdej
formuły z parametrami diagram zawiera odpowiednie
zdanie lub jego negację. Semantyka podstawienia
(twierdzenie \ref{thm:sem-substitution}) daje zachowanie
i odbijanie prawdziwości formuły przy obrazach parametrów,
czyli elementarność.
Odwrotna implikacja wynika z zachowania formuł bez
kwantyfikatorów przez zanurzenie lub wszystkich
formuł przez zanurzenie elementarne.
\end{proof}

\begin{lemat}[Przedłużenie zanurzenia do rozszerzenia elementarnego]
\label{lem:mod-mc-robinson-diagram}
Dla $\mathfrak M\subseteq\mathfrak N$ równoważne są warunki:
\begin{enumerate}[label=(\roman*)]
\item $\mathfrak M\preccurlyeq_{\exists}\mathfrak N$;
\item istnieje rozszerzenie elementarne
$\mathfrak M'\succeq\mathfrak M$ i zanurzenie
$f:\mathfrak N\to\mathfrak M'$ ustalające każdy element $M$.
\end{enumerate}
\end{lemat}
\begin{proof}
Załóżmy (i). Nazwy elementów $M$ w obu diagramach
wybieramy identycznie i rozważamy teorię
\[
\operatorname{ED}(\mathfrak M)
\cup\operatorname{Diag}(\mathfrak N).
\]
Każdy jej skończony fragment ma model: wybrane zdania
z drugiego diagramu tworzą koniunkcję bez kwantyfikatorów
$\theta(\bar c,\bar d)$, gdzie $\bar c$ są nazwami
elementów $M$, a $\bar d$ nazwami skończenie wielu
elementów $N\setminus M$. Zastępując nazwy $\bar d$
świeżymi zmiennymi $\bar y$, otrzymujemy formułę
$\exists\bar y\,\theta(\bar c,\bar y)$ prawdziwą
w $\mathfrak N$. Domknięcie egzystencjalne przenosi
ją na $\mathfrak M$. Interpretujemy zatem nazwy
$\bar d$ odpowiednimi świadkami w $M$, a pozostałe
nowe nazwy dowolnymi elementami $M$.
Przy nazwach z $M$ interpretowanych standardowo
wszystkie zdania diagramu elementarnego pozostają
prawdziwe. Zwartość (twierdzenie \ref{thm:compactness})
daje model całej teorii.

Na mocy lematu \ref{lem:mod-mc-diagram-embedding}
uzyskujemy elementarne zanurzenie $\mathfrak M$ i
zanurzenie $\mathfrak N$ w ten sam redukt, zgodne na $M$.
Przenosząc strukturę reduktu wzdłuż bijekcji, która
utożsamia obraz $a\in M$ z samym $a$, a pozostałe
elementy wysyła na świeże elementy spoza $M$, uzyskujemy
rzeczywiste rozszerzenie elementarne $\mathfrak M'$
oraz wymagane $f$. Świeży zbiór odpowiedniej mocy
można wybrać jako zbiór liczb porządkowych większych
od rangi $M$.

Odwrotnie, jeśli formuła egzystencjalna z parametrami
z $M$ jest prawdziwa w $\mathfrak N$, jej prawdziwość
przenosi się przez $f$ do $\mathfrak M'$ na mocy
lematu \ref{lem:mod-mc-preservation} i niezmienniczości
względem izomorfizmów. Parametry pozostają niezmienione.
Elementarność $\mathfrak M\preccurlyeq\mathfrak M'$
przenosi formułę do $\mathfrak M$, co daje (i).
\end{proof}

\begin{twierdzenie}[Kryterium Robinsona]
\label{thm:mod-mc-robinson}
Dla spełnialnej teorii $T$ równoważne są warunki:
\begin{enumerate}[label=(\roman*)]
\item $T$ jest modelowo zupełna;
\item każdy model $T$ jest egzystencjalnie domkniętym
modelem $T$;
\item dla dowolnych modeli $\mathfrak M\subseteq\mathfrak N$
teorii $T$ zachodzi $\mathfrak M\preccurlyeq_{\exists}\mathfrak N$.
\end{enumerate}
\end{twierdzenie}
\begin{proof}
Implikacja (i)$\Rightarrow$(ii) wynika z definicji
elementarności, a równoważność (ii) i (iii) jest
bezpośrednim przepisaniem kwantyfikatorów w definicji
domknięcia egzystencjalnego.

Załóżmy (iii) i ustalmy modele
$\mathfrak M_0=\mathfrak M\subseteq\mathfrak N_0=\mathfrak N$.
Z lematu \ref{lem:mod-mc-robinson-diagram} otrzymujemy
rozszerzenie elementarne $\mathfrak M_1$ struktury
$\mathfrak M_0$ oraz zanurzenie $\mathfrak N_0$ w
$\mathfrak M_1$ ustalające $M_0$. Przenosząc strukturę
$\mathfrak M_1$ na izomorficzną kopię, możemy założyć,
że $\mathfrak N_0\subseteq\mathfrak M_1$, przy zachowaniu
$\mathfrak M_0\preccurlyeq\mathfrak M_1$.
Ponieważ $\mathfrak M_1$ jest modelem $T$, ponownie
stosujemy (iii), tym razem do
$\mathfrak N_0\subseteq\mathfrak M_1$, i konstruujemy
$\mathfrak N_1\succeq\mathfrak N_0$ zawierającą
$\mathfrak M_1$. Powtarzamy naprzemiennie konstrukcję:
\[
\mathfrak M_0\subseteq\mathfrak N_0
\subseteq\mathfrak M_1\subseteq\mathfrak N_1
\subseteq\mathfrak M_2\subseteq\cdots,
\]
przy czym dla każdego $i\in\N$
\[
\mathfrak M_i\preccurlyeq\mathfrak M_{i+1},
\qquad
\mathfrak N_i\preccurlyeq\mathfrak N_{i+1}.
\]
Wybory kolejnych rozszerzeń można ograniczyć do
zbioru: na mocy dolnego twierdzenia Löwenheima--Skolema
każde rozszerzenie elementarne w konstrukcji można
zmniejszyć do podstruktury elementarnej zawierającej
obraz większego z dotychczasowych modeli i o mocy
nie większej niż
$\mu=\max(\aleph_0,|\mathcal L|,|N_0|)$
(twierdzenie \ref{thm:downward-ls}).
Jeżeli $\mathfrak M_i$ jest elementarna w dużym
rozszerzeniu, a wybrana podstruktura elementarna
tego rozszerzenia zawiera $\mathfrak M_i$, to
$\mathfrak M_i$ jest elementarna także w wybranej
podstrukturze: obie zgadzają się z dużym modelem
co do formuł z parametrami z $M_i$.
Tak samo postępujemy w kroku dla $\mathfrak N_i$.
Przenosimy nowe elementy do kolejnych rozłącznych
zbiorów mocy $\mu$, rozłącznych też z $N_0$.
Wtedy aksjomat wyboru pozwala ustalić wybór rozszerzenia
na każdym etapie w obrębie zbioru struktur
na podzbiorach ustalonej sumy tych zbiorów.

Oba ciągi mają tę samą strukturę sumy $\mathfrak U$,
ponieważ przeplatają się i ich interpretacje są zgodne
na wcześniejszych uniwersach. Z lematu o sumach łańcuchów
elementarnych \ref{lem:mod-saturation-chain} wynika
\[
\mathfrak M_0\preccurlyeq\mathfrak U,
\qquad \mathfrak N_0\preccurlyeq\mathfrak U.
\]
Dla dowolnej formuły z parametrami z $M_0$ jej
prawdziwość w $\mathfrak M_0$ i $\mathfrak N_0$
jest więc równoważna prawdziwości w $\mathfrak U$.
Stąd $\mathfrak M_0\preccurlyeq\mathfrak N_0$,
co dowodzi (i).
\end{proof}

\begin{figure}[htbp]
\centering
\begin{tikzpicture}[>=Stealth,x=1cm,y=1cm]
\foreach \x/\i in {0/0,3/1,6/2}{
 \node[draw,rounded corners,minimum width=1.1cm] (m\i) at (\x,1.8) {$\mathfrak M_{\i}$};
 \node[draw,rounded corners,minimum width=1.1cm] (n\i) at (\x+1,0) {$\mathfrak N_{\i}$};
 \draw[->,gray] (m\i)--node[left,inner sep=1pt] {$\subseteq$}(n\i);
}
\foreach \i/\j in {0/1,1/2}{
 \draw[->,blue!65!black,thick] (m\i)--node[above] {$\preccurlyeq$}(m\j);
 \draw[->,blue!65!black,thick] (n\i)--node[below] {$\preccurlyeq$}(n\j);
 \draw[->,gray] (n\i)--node[right,inner sep=1pt] {$\subseteq$}(m\j);
}
\node[draw,rounded corners] (u) at (9,0.9) {$\mathfrak U$};
\draw[->,blue!65!black,dashed] (m2)--(u);
\draw[->,blue!65!black,dashed] (n2)--(u);
\end{tikzpicture}
\caption{Przeplatające się łańcuchy w dowodzie kryterium
Robinsona. Niebieskie strzałki oznaczają elementarność.}
\label{fig:mod-mc-robinson-chain}
\end{figure}

Rysunek \ref{fig:mod-mc-robinson-chain} przedstawia
konstrukcję z dowodu twierdzenia \ref{thm:mod-mc-robinson}.
Istotne jest to, że samo domknięcie egzystencjalne jednej
podstruktury nie musi oznaczać jej elementarności.
Kryterium dotyczy wszystkich par modeli teorii, dzięki
czemu konstrukcję można powtarzać.

\begin{przyklad}[Domknięcie egzystencjalne bez elementarności]
\label{ex:mod-mc-ec-not-elementary}
W języku $\{<\}$ niech
$\mathfrak M=(\Q,<)$, a $\mathfrak N$ będzie porządkiem
otrzymanym przez dodanie jednego nowego największego
elementu $t$ do $\Q$.
Wtedy $\mathfrak M\preccurlyeq_{\exists}\mathfrak N$,
ale $\mathfrak M\not\preccurlyeq\mathfrak N$.
Co więcej, w $\mathfrak M$ i $\mathfrak N$ zachodzą
dokładnie te same zdania egzystencjalne.
\end{przyklad}
\begin{proof}
Niech formuła egzystencjalna
$\exists\bar y\,\theta(\bar a,\bar y)$ z parametrami
$\bar a$ z $\Q$ zachodzi w $\mathfrak N$.
Wybieramy krotkę świadków $\bar b$.
Jeżeli żadna współrzędna $\bar b$ nie jest równa $t$,
to wszystkie parametry i świadkowie należą do $\Q$,
więc macierz zachodzi również w $\mathfrak M$.
W przeciwnym razie wybieramy liczbę wymierną $q$
większą od wszystkich wymiernych współrzędnych
$\bar a,\bar b$ (jeżeli ich nie ma, wybieramy $q=0$).
Zastępujemy każdą współrzędną równą $t$ przez $q$.
W ten sposób zachowujemy wszystkie równości
i porównania w połączonej krotce, a więc również
prawdziwość macierzy bez kwantyfikatorów.
Otrzymujemy świadków w $\mathfrak M$.
Z definicji zachodzi domknięcie egzystencjalne;
w szczególności oba porządki zgadzają się co do
zdań egzystencjalnych, przez lemat
\ref{lem:mod-mc-preservation}.
Jednak zdanie
$\exists x\,\forall y\,\neg(x<y)$
jest prawdziwe w $\mathfrak N$, ze świadkiem $t$,
a fałszywe w $\mathfrak M$, ponieważ dla każdego
$x\in\Q$ mamy $x<x+1$.
Inkluzja nie jest więc elementarna.
\end{proof}

\section{Charakterystyka przez równoważniki egzystencjalne}
\begin{lemat}[Diagram a typ egzystencjalny]
\label{lem:mod-mc-existential-type}
Niech $\mathfrak M,\mathfrak N$ będą strukturami
tego samego języka, a $\bar a,\bar b$ skończonymi
krotkami tej samej długości.
Jeżeli każda formuła egzystencjalna prawdziwa przy
$\bar a$ w $\mathfrak M$ jest prawdziwa przy $\bar b$
w $\mathfrak N$, to istnieje rozszerzenie elementarne
$\mathfrak N'\succeq\mathfrak N$ i zanurzenie
$f:\mathfrak M\to\mathfrak N'$ wysyłające $\bar a$
na $\bar b$. Odwrotnie, istnienie takiego zanurzenia
implikuje wskazane zachowanie formuł egzystencjalnych.
Dopuszczamy również krotki puste.
\end{lemat}
\begin{proof}
Najpierw zauważmy, że obie krotki mają ten sam układ
równości, ponieważ równości i ich negacje są formułami
egzystencjalnymi z pustą listą kwantyfikatorów.
Możemy więc spójnie utożsamić nazwy odpowiadających
sobie współrzędnych w diagramach
$\operatorname{Diag}(\mathfrak M)$ i
$\operatorname{ED}(\mathfrak N)$, pozostawiając wszystkie
pozostałe nazwy rozłączne.

Skończony fragment diagramu $\mathfrak M$ jest
koniunkcją bez kwantyfikatorów dotyczącą nazw
współrzędnych $\bar a$ i skończenie wielu innych
elementów. Zastępujemy te inne nazwy świeżymi zmiennymi
i wiążemy je egzystencjalnie. Otrzymana formuła
egzystencjalna jest prawdziwa przy $\bar a$,
więc z założenia także przy $\bar b$ w $\mathfrak N$.
Jej świadkowie pozwalają zinterpretować dodatkowe
nazwy diagramu $\mathfrak M$ w $\mathfrak N$,
nie zmieniając interpretacji nazw z diagramu
elementarnego $\mathfrak N$.
W konsekwencji każdy skończony fragment sumy diagramów
ma model. Zwartość daje model całej sumy, a lemat
\ref{lem:mod-mc-diagram-embedding} daje dwa zanurzenia
(z których zanurzenie $\mathfrak N$ jest elementarne)
zgodne na wskazanych nazwach.
Jak w dowodzie lematu \ref{lem:mod-mc-robinson-diagram}
przenosimy wspólną strukturę na izomorficzną kopię
zawierającą rzeczywiście $\mathfrak N$ jako podstrukturę
elementarną.

Odwrotnie, jeśli istnieje takie $f$, to formuła
egzystencjalna prawdziwa przy $\bar a$ w $\mathfrak M$
jest prawdziwa przy $f(\bar a)=\bar b$ w $\mathfrak N'$
przez lemat \ref{lem:mod-mc-preservation}, zastosowany
do obrazu zanurzenia. Elementarność
$\mathfrak N\preccurlyeq\mathfrak N'$ przenosi jej
prawdziwość do $\mathfrak N$. W przypadku pustych
krotek w diagramach nie utożsamiamy żadnych nowych
nazw i argument daje ten sam wniosek dla zdań.
\end{proof}

\begin{twierdzenie}[Zupełność modelowa a formuły egzystencjalne]
\label{thm:mod-mc-existential-formulas}
Dla spełnialnej teorii $T$ równoważne są warunki:
\begin{enumerate}[label=(\roman*)]
\item $T$ jest modelowo zupełna;
\item każda formuła jest równoważna modulo $T$ formule
egzystencjalnej bez dodatkowych zmiennych wolnych;
\item każda formuła jest równoważna modulo $T$ formule
uniwersalnej bez dodatkowych zmiennych wolnych.
\end{enumerate}
\end{twierdzenie}
\begin{proof}
Załóżmy (i), ustalmy formułę $\varphi(\bar x)$
i jej realizację $\bar a$ w pewnym modelu
$\mathfrak M\models T$. Niech $p$ będzie zbiorem
wszystkich formuł egzystencjalnych w liście zmiennych
$\bar x$, prawdziwych przy $\bar a$.
Twierdzimy, że $p$ wymusza $\varphi$ modulo $T$.
Istotnie, gdyby $\bar b$ w pewnym modelu
$\mathfrak N\models T$ realizowała $p$, ale nie
$\varphi$, to lemat \ref{lem:mod-mc-existential-type}
dałby zanurzenie $\mathfrak M$ w rozszerzenie
elementarne $\mathfrak N'$ modelu $\mathfrak N$,
wysyłające $\bar a$ na $\bar b$.
Oba modele spełniają $T$, więc z (i) zanurzenie
jest elementarne. Prawdziwość $\varphi$ przenosi się
na $\bar b$ w $\mathfrak N'$ i następnie w $\mathfrak N$,
co jest sprzecznością.

Po dodaniu świeżych nazw $\bar c$ dla wolnych zmiennych
teoria
\[
T\cup\{\eta(\bar c):\eta\in p\}
\cup\{\neg\varphi(\bar c)\}
\]
jest niespełnialna. Zwartość daje skończony fragment
$p_0\subseteq p$ wymuszający $\varphi$ modulo $T$.
Koniunkcja formuł z $p_0$ jest równoważna formule
egzystencjalnej: przemianowujemy związane zmienne
w członach na parami rozłączne listy, a następnie
przenosimy wszystkie kwantyfikatory egzystencjalne
przed koniunkcję macierzy. Równoważność wynika z tego,
że świadków dla poszczególnych członów można wybrać
niezależnie przy tym samym wartościowaniu zmiennych
wolnych. Dla $p_0=\varnothing$ przyjmujemy $\top$.
Uzyskana formuła egzystencjalna jest prawdziwa
przy $\bar a$ i implikuje $\varphi$ modulo $T$.

Niech $\Theta$ będzie zbiorem wszystkich formuł
egzystencjalnych w liście $\bar x$ implikujących
$\varphi$ modulo $T$. Każda realizacja $\varphi$
w modelu $T$ spełnia pewną formułę z $\Theta$,
co właśnie wykazaliśmy. Zatem teoria
\[
T\cup\{\varphi(\bar c)\}
\cup\{\neg\theta(\bar c):\theta\in\Theta\}
\]
jest niespełnialna. Zwartość daje skończoną alternatywę
formuł z $\Theta$ równoważną $\varphi$ modulo $T$.
Taka alternatywa jest równoważna formule egzystencjalnej:
po przemianowaniu zmiennych przenosimy wszystkie
kwantyfikatory przed alternatywę macierzy.
Jeśli jeden człon jest prawdziwy, dla jego zmiennych
wybieramy świadków, a dla pozostałych dowolne
elementy niepustego uniwersum. Implikacja odwrotna
wynika z prawdziwości co najmniej jednego członu
macierzy. Pusta alternatywa to $\bot$.
Otrzymujemy (ii), również w przypadku zdań.

Warunki (ii) i (iii) są równoważne: stosujemy odpowiednią
własność do negacji formuły, a następnie przenosimy
negację przez blok kwantyfikatorów, zamieniając
egzystencjalne na uniwersalne lub odwrotnie.

Na koniec załóżmy (ii) i niech
$\mathfrak M\subseteq\mathfrak N$ będą modelami $T$.
Jeśli $\mathfrak M\models\varphi(\bar a)$, stosujemy
egzystencjalny równoważnik $\varphi$ i lemat
\ref{lem:mod-mc-preservation}, otrzymując
$\mathfrak N\models\varphi(\bar a)$.
Jeśli $\mathfrak M\not\models\varphi(\bar a)$,
stosujemy ten sam argument do $\neg\varphi$,
otrzymując $\mathfrak N\not\models\varphi(\bar a)$.
Zatem inkluzja jest elementarna i zachodzi (i).
\end{proof}

\begin{przyklad}[Zupełność modelowa bez eliminacji kwantyfikatorów]
\label{ex:mod-mc-no-qe}
Niech $T$ będzie zupełną teorią struktury $(\R,0,1,+,-,\cdot)$
w języku pierścieni bez symbolu porządku.
W każdym modelu $T$ relacja
\[
a<b\quad\Longleftrightarrow\quad
\exists z\,(z\ne0\land z^2=b-a)
\]
jest jedynym rozszerzeniem o porządek, które daje
model teorii ciał rzeczywiście domkniętych.
Aby to uzasadnić, tłumaczymy wszystkie zdania
języka uporządkowanych pierścieni przez podstawienie
powyższej definicji za atomy porządku, używając
świeżych zmiennych związanych. Indukcja po formułach
daje zgodność prawdziwości z rozwinięciem struktury.
W $\R$ definicja jest zgodna ze zwykłym porządkiem,
zatem tłumaczenia aksjomatów ciał rzeczywiście
domkniętych (lemat \ref{lem:mod-qe-rcf-axioms})
należą do $T$. W każdym modelu $T$ zdefiniowana
relacja daje więc ciało rzeczywiście domknięte.
Jedyność porządku wynika z przykładu
\ref{ex:mod-qe-rcf-square}: dodatniość jest tam
równoważna byciu niezerowym kwadratem.

Każde zanurzenie modeli $T$ zachowuje dodatniość,
bo zachowuje niezerowe kwadraty. Odbija ją również:
jeżeli $a\ne0$ nie jest dodatni, to $-a$ jest dodatni,
więc jego obraz jest dodatni, a obraz $a$ ujemny.
Element zerowy nie jest dodatni w żadnym z modeli.
Zanurzenie jest zatem zanurzeniem uporządkowanych ciał.
Eliminacja kwantyfikatorów dla ciał rzeczywiście
domkniętych (twierdzenie \ref{thm:mod-qe-rcf})
daje elementarność w języku z porządkiem,
a więc także w redukcie. Teoria $T$ jest modelowo zupełna.

Nie ma ona jednak eliminacji kwantyfikatorów w języku
pierścieni. W standardowym modelu $\R$ weźmy liczbę
przestępną $t>0$. Zarówno $t$, jak i $-t$ są przestępne
nad $\Q$: gdyby niezerowy wielomian o współczynnikach
wymiernych zerował się w $-t$, to po zmianie znaku
zmiennej otrzymalibyśmy niezerowy wielomian zerujący
się w $t$. Istnienie dodatniej liczby przestępnej
wynika z nieprzeliczalności $\R$ oraz przeliczalności
zbioru liczb algebraicznych (jest sumą skończonych
zbiorów pierwiastków przeliczalnej rodziny niezerowych
wielomianów wymiernych).
Każdy atom w jednej zmiennej jest równością wielomianową
o współczynnikach całkowitych; w obu punktach
zachodzi dokładnie wtedy, gdy różnica wielomianów
jest wielomianem zerowym. Krotki $(t)$ i $(-t)$
mają więc ten sam typ bez kwantyfikatorów.
Rozróżnia je formuła $\exists z\,(z^2=x)$, prawdziwa
przy $t$ i fałszywa przy $-t$.
Z kryterium \ref{thm:mod-qe-semantic} wynika brak
eliminacji kwantyfikatorów w języku pierścieni.
\end{przyklad}

\begin{wniosek}[Test dla zupełności modelowej]
\label{cor:mod-mc-neg-existential}
Aby wykazać zupełność modelową spełnialnej teorii $T$,
wystarczy wykazać, że negacja każdej formuły
egzystencjalnej jest równoważna modulo $T$ formule
egzystencjalnej bez dodatkowych zmiennych wolnych.
Warunek ten jest również konieczny.
\end{wniosek}
\begin{proof}
Konieczność wynika z twierdzenia
\ref{thm:mod-mc-existential-formulas}.
Dla dostateczności stosujemy indukcję po budowie
formuły. Atomy są już egzystencjalne, z pustym
blokiem kwantyfikatorów. Koniunkcje i alternatywy
egzystencjalnych równoważników zastępujemy
egzystencjalnymi formułami w postaci opisanej
w dowodzie twierdzenia
\ref{thm:mod-mc-existential-formulas}.
Negację obsługujemy za pomocą założenia.
Kwantyfikator egzystencjalny dołączamy do bloku
egzystencjalnego po ewentualnej zmianie nazw
zmiennych związanych. Kwantyfikator uniwersalny
wyrażamy przez negację, kwantyfikator egzystencjalny
i negację. Pozostałe klasyczne spójniki wyrażamy
przez negację, koniunkcję i alternatywę.
Wszystkie kroki zachowują równoważność modulo $T$
i nie wprowadzają nowych zmiennych wolnych.
Każda formuła ma zatem równoważnik egzystencjalny,
więc stosujemy to samo twierdzenie.
\end{proof}

\section{Teorie uniwersalne}
\begin{definicja}[Konsekwencje uniwersalne]
\label{def:mod-mc-universal-theory}
Dla teorii $T$ oznaczamy przez $T_{\forall}$ zbiór
wszystkich zdań uniwersalnych będących konsekwencjami
semantycznymi $T$. Teorię nazywamy uniwersalną,
jeżeli można ją aksjomatyzować zdaniami uniwersalnymi,
czyli ma ona te same modele co pewien zbiór takich zdań.
Równoważnie, $T$ i $T_{\forall}$ mają te same modele.
\end{definicja}

Ostatnia równoważność wynika z tego, że każdy
aksjomat uniwersalny dowolnej teorii równoważnej
semantycznie $T$ należy do $T_{\forall}$,
a każdy model $T$ spełnia wszystkie jej konsekwencje.
Jeżeli odwrotnie $T$ oraz $T_{\forall}$ mają te same
modele, to sam zbiór $T_{\forall}$ stanowi poszukiwaną
aksjomatyzację.
Rozróżniamy syntaktyczną postać aksjomatów od
aksjomatyzowalności: teoria uniwersalna może zawierać
także zdania o innej postaci syntaktycznej.

\begin{twierdzenie}[Zanurzanie w modele teorii]
\label{thm:mod-mc-universal-embedding}
Niech $T$ będzie spełnialną teorią i $\mathfrak A$
strukturą jej języka. Następujące warunki są równoważne:
\begin{enumerate}[label=(\roman*)]
\item $\mathfrak A\models T_{\forall}$;
\item istnieje model $\mathfrak B\models T$ zawierający
$\mathfrak A$ jako podstrukturę.
\end{enumerate}
\end{twierdzenie}
\begin{proof}
Jeżeli zachodzi (ii), to każde zdanie z $T_{\forall}$
jest prawdziwe w $\mathfrak B$, a więc również
w $\mathfrak A$ na mocy twierdzenia
\ref{thm:model-universal-substructure}.

Załóżmy (i). Wykażemy spełnialność teorii
\[
T\cup\operatorname{Diag}(\mathfrak A).
\]
Gdyby nie miała modelu, zwartość dałaby skończony
fragment diagramu, którego koniunkcję oznaczamy
przez $\theta(c_{a_1},\ldots,c_{a_n})$, sprzeczną z $T$.
Wszystkie występujące świeże nazwy elementów
$\mathfrak A$ zastępujemy parami różnymi zmiennymi
$\bar x$, uzyskując formułę bez kwantyfikatorów
$\theta(\bar x)$ języka pierwotnego.
Niespełnialność oznacza
\[
T\models\forall\bar x\,\neg\theta(\bar x):
\]
gdyby w modelu $T$ istniała krotka spełniająca
$\theta$, można by zinterpretować nowe nazwy
jej współrzędnymi i otrzymać model $T$ wraz
z wybranym fragmentem diagramu. Nazwy niewystępujące
w tym fragmencie interpretujemy dowolnie.
Powyższe zdanie należy do $T_{\forall}$, ale jest
fałszywe w $\mathfrak A$, ponieważ
$\mathfrak A\models\theta(\bar a)$. Sprzeczność.

Zatem suma teorii i diagramu ma model. Lemat
\ref{lem:mod-mc-diagram-embedding} daje zanurzenie
$\mathfrak A$ w redukt tego modelu. Przenosząc
redukt na izomorficzną kopię utożsamiającą
obraz $a$ z $a$ i używając świeżych elementów
dla reszty uniwersum, otrzymujemy (ii).
Argument obejmuje przypadek $n=0$, kiedy
nie dodajemy kwantyfikatorów.
\end{proof}

\begin{twierdzenie}[Twierdzenie o zachowaniu przy podstrukturach]
\label{thm:mod-mc-universal-preservation}
Spełnialna teoria $T$ jest uniwersalna dokładnie wtedy,
gdy każda niepusta podstruktura dowolnego modelu $T$
jest modelem $T$.
\end{twierdzenie}
\begin{proof}
Jeżeli $T$ jest uniwersalna, jej aksjomaty zachowują
prawdziwość przy przejściu do podstruktury na mocy
twierdzenia \ref{thm:model-universal-substructure}.
Odwrotnie, niech klasa modeli $T$ będzie domknięta
na podstruktury. Każdy model $T_{\forall}$ jest
podstrukturą pewnego modelu $T$ na mocy twierdzenia
\ref{thm:mod-mc-universal-embedding}, a więc
sam spełnia $T$. Ponieważ każdy model $T$ spełnia
$T_{\forall}$, obie teorie mają te same modele.
\end{proof}

\begin{przyklad}[Ciała a dziedziny całkowitości]
\label{ex:mod-mc-fields-universal}
W języku pierścieni z jedynką teoria ciał nie jest
uniwersalna. Pierścień $\mathbb Z$ jest podstrukturą
ciała $\Q$, ale nie jest ciałem: element $2$ nie ma
odwrotności w $\mathbb Z$. Z twierdzenia
\ref{thm:mod-mc-universal-preservation} wynika brak
aksjomatyzacji uniwersalnej.

Klasa modeli uniwersalnych konsekwencji teorii ciał
składa się dokładnie z dziedzin całkowitości.
Istotnie, podstruktura ciała w tym języku jest
pierścieniem przemiennym z $1\ne0$ bez dzielników
zera. Odwrotnie, każda dziedzina całkowitości $D$
zanurza się w ciało swoich ułamków: jego elementy
reprezentujemy przez pary $(a,b)$ z $a,b\in D$, $b\ne0$,
identyfikując $(a,b)$ z $(c,d)$, gdy $ad=bc$.
Zanurzenie jest dane przez $a\mapsto[(a,1)]$,
a niezerowy ułamek $[(a,b)]$ ma odwrotność $[(b,a)]$.
Dla pełności przypomnijmy uzasadnienie tej konstrukcji.
Przechodniość relacji równoważności wynika z tego,
że z $ad=bc$ i $cf=de$ otrzymujemy $adf=bde$,
a po skróceniu przez niezerowe $d$ mamy $af=be$.
Zwrotność i symetria są bezpośrednie.
Działania definiujemy wzorami
\[
[(a,b)]+[(c,d)]=[(ad+bc,bd)],
\qquad
[(a,b)]\,[(c,d)]=[(ac,bd)].
\]
Mianowniki są niezerowe. Jeśli $ab'=a'b$ oraz
$cd'=c'd$, to
\[
\begin{aligned}
(ad+bc)b'd'&=(a'd'+b'c')bd,\\
acb'd'&=a'c'bd,
\end{aligned}
\]
co wynika z rozdzielności i podstawienia obu równości.
Działania są zatem dobrze określone.
Aksjomaty pierścienia wynikają z działań w $D$:
dla dowolnej skończonej listy ułamków można
sprowadzić je do wspólnego niezerowego mianownika;
po rozwinięciu działań obie strony odpowiedniej
tożsamości pierścieniowej mają, po pomnożeniu
przez wspólny mianownik, równe liczniki na mocy
aksjomatów $D$. Zero i jedynka to odpowiednio
$[(0,1)]$ i $[(1,1)]$, różne, gdyż $0\ne1$ w $D$.
Element $[(a,b)]$ jest zerem dokładnie wtedy,
gdy $a=0$, a gdy $a\ne0$, wskazana odwrotność
jest dobrze określona i ich iloczyn to
$[(ab,ba)]=[(1,1)]$.
Zanurzenie jest różnowartościowe, ponieważ
$[(a,1)]=[(c,1)]$ implikuje $a=c$,
oraz zachowuje stałe i działania przez podane wzory.
Stosujemy twierdzenie \ref{thm:mod-mc-universal-embedding}.
\end{przyklad}

\section{Teorie uniwersalno-egzystencjalne}
\begin{definicja}[Zdania uniwersalno-egzystencjalne]
\label{def:mod-mc-forall-exists}
Zdaniem uniwersalno-egzystencjalnym nazywamy zdanie
postaci
\[
\forall\bar x\,\exists\bar y\,\theta(\bar x,\bar y),
\]
gdzie macierz $\theta$ nie zawiera kwantyfikatorów,
a listy zmiennych są skończone, rozłączne, bez powtórzeń
i obejmują wszystkie zmienne wolne macierzy.
Każda z list może być pusta. Zbiór wszystkich takich konsekwencji teorii
$T$ oznaczamy przez $T_{\forall\exists}$.
Teorię nazywamy uniwersalno-egzystencjalną, jeżeli
ma aksjomatyzację zdaniami tej postaci, czyli ma
te same modele co $T_{\forall\exists}$.
\end{definicja}

\begin{lemat}[Zachowanie przy sumach łańcuchów]
\label{lem:mod-mc-chain-preservation}
Niech $(\mathfrak M_i)_{i<\delta}$ będzie niepustym
rosnącym łańcuchem podstruktur, indeksowanym liczbą
porządkową $\delta>0$. Jeżeli każde $\mathfrak M_i$
spełnia zdanie uniwersalno-egzystencjalne $\sigma$,
to spełnia je również struktura sumy
$\mathfrak U=\bigcup_{i<\delta}\mathfrak M_i$.
\end{lemat}
\begin{proof}
Na sumie uniwersów interpretujemy stałe i funkcje
zgodnie z interpretacjami w członach łańcucha.
Każda skończona krotka elementów sumy należy
do pewnego wspólnego członu, ponieważ maksimum
skończonego zbioru indeksów jest indeksem łańcucha.
Wartości funkcji i relacji na tej krotce nie zależą
od wyboru członu, dzięki inkluzjom podstruktur.
Dla symboli o arności zero wybieramy dowolny człon;
ich interpretacje są zgodne. Otrzymujemy strukturę
$\mathfrak U$ zawierającą każdy człon jako podstrukturę.

Niech $\sigma=\forall\bar x\,\exists\bar y\,\theta$.
Dla dowolnej krotki $\bar a$ z sumy wybieramy człon
$\mathfrak M_i$ zawierający wszystkie jej współrzędne
(gdy krotka jest pusta, dowolny człon). Ponieważ
$\mathfrak M_i\models\sigma$, istnieje krotka
$\bar b$ w tym członie spełniająca
$\theta(\bar a,\bar b)$. Prawdziwość macierzy bez
kwantyfikatorów zachowuje się w strukturze sumy
na mocy definicji podstruktury, jak w lemacie
\ref{lem:mod-mc-preservation}. Zatem $\bar b$ jest
świadkiem w sumie. Ponieważ $\bar a$ była dowolna,
otrzymujemy $\mathfrak U\models\sigma$.
\end{proof}

\begin{lemat}[Rozszerzenie spełniające teorię]
\label{lem:mod-mc-forall-exists-extension}
Dla spełnialnej teorii $T$ oraz struktury
$\mathfrak A$ jej języka równoważne są warunki:
\begin{enumerate}[label=(\roman*)]
\item $\mathfrak A\models T_{\forall\exists}$;
\item istnieje $\mathfrak B\models T$ taka, że
$\mathfrak A\preccurlyeq_{\exists}\mathfrak B$.
\end{enumerate}
\end{lemat}
\begin{proof}
Jeżeli zachodzi (ii), to dla dowolnego zdania
$\forall\bar x\,\exists\bar y\,\theta$ z
$T_{\forall\exists}$ i dowolnej krotki $\bar a$ z $A$
mamy $\mathfrak B\models\exists\bar y\,\theta(\bar a,\bar y)$.
Z domknięcia egzystencjalnego ta formuła zachodzi
w $\mathfrak A$, więc $\mathfrak A$ spełnia całe zdanie.
Dotyczy to także pustej krotki i pustego bloku
egzystencjalnego. Otrzymujemy (i).

Załóżmy (i). Niech $\operatorname{UD}(\mathfrak A)$ oznacza
zbiór wszystkich zdań uniwersalnych prawdziwych
w strukturze $\mathfrak A$ z nazwami wszystkich
jej elementów. Zawiera on cały diagram bez
kwantyfikatorów, ponieważ dopuszczamy pusty
blok uniwersalny.

Twierdzimy, że $T\cup\operatorname{UD}(\mathfrak A)$
jest spełnialna. Gdyby nie, zwartość dałaby skończony
fragment diagramu uniwersalnego sprzeczny z $T$.
Po przemianowaniu zmiennych związanych na rozłączne
listy koniunkcję tego fragmentu można zapisać jako
\[
\forall\bar y\,\theta(\bar c,\bar y),
\]
gdzie $\theta$ nie zawiera kwantyfikatorów.
Jest to równoważność semantyczna: w niepustym
uniwersum jednoczesna uniwersalna prawdziwość
każdej z macierzy jest równoważna prawdziwości
ich koniunkcji przy wszystkich wartościach
połączonej listy zmiennych.
Nazwy $\bar c$ są dokładnie świeżymi stałymi
występującymi w skończonym fragmencie; zastępujemy je
świeżymi zmiennymi $\bar x$, rozłącznymi z listą
$\bar y$. Niespełnialność oznacza
\[
T\models\forall\bar x\,\exists\bar y\,
\neg\theta(\bar x,\bar y).
\]
Rzeczywiście, gdyby jakiś model $T$ miał krotkę
$\bar a$ z $\forall\bar y\,\theta(\bar a,\bar y)$,
interpretacja świeżych nazw przez tę krotkę
spełniałaby fragment diagramu wraz z $T$.
Otrzymane zdanie uniwersalno-egzystencjalne należy
do $T_{\forall\exists}$, lecz jest fałszywe
w $\mathfrak A$ przy krotce rzeczywistych wartości
nazw $\bar c$. Sprzeczność.

Model sumy teorii i diagramu uniwersalnego zawiera,
po przeniesieniu struktury wzdłuż izomorfizmu,
$\mathfrak A$ jako podstrukturę, na mocy lematu
\ref{lem:mod-mc-diagram-embedding}.
Niech $\mathfrak B$ będzie tak uzyskanym reduktem.
Każda formuła uniwersalna z parametrami z $A$
prawdziwa w $\mathfrak A$ jest prawdziwa w $\mathfrak B$,
ponieważ odpowiednie zdanie z nazwami należy
do diagramu uniwersalnego. Jeżeli więc formuła
egzystencjalna z parametrami z $A$ byłaby prawdziwa
w $\mathfrak B$ i fałszywa w $\mathfrak A$, jej
negacja miałaby uniwersalny równoważnik prawdziwy
w $\mathfrak A$, a zatem również w $\mathfrak B$,
co daje sprzeczność. To dowodzi
$\mathfrak A\preccurlyeq_{\exists}\mathfrak B$.
\end{proof}

\begin{twierdzenie}[Twierdzenie Changa--Łosia--Suszki]
\label{thm:mod-mc-inductive}
Dla spełnialnej teorii $T$ równoważne są warunki:
\begin{enumerate}[label=(\roman*)]
\item $T$ jest uniwersalno-egzystencjalna;
\item suma każdego niepustego rosnącego łańcucha
modeli $T$ jest modelem $T$;
\item suma każdego rosnącego ciągu modeli $T$
indeksowanego przez $\N$ jest modelem $T$.
\end{enumerate}
\end{twierdzenie}
\begin{proof}
Implikacja (i)$\Rightarrow$(ii) wynika z lematu
\ref{lem:mod-mc-chain-preservation} zastosowanego
do każdego aksjomatu. Implikacja (ii)$\Rightarrow$(iii)
jest natychmiastowa.

Załóżmy (iii) i niech
$\mathfrak A_0\models T_{\forall\exists}$.
Z lematu \ref{lem:mod-mc-forall-exists-extension}
bierzemy $\mathfrak B_0\models T$ z
$\mathfrak A_0\preccurlyeq_{\exists}\mathfrak B_0$.
Lemat \ref{lem:mod-mc-robinson-diagram} pozwala
zanurzyć $\mathfrak B_0$ nad $\mathfrak A_0$
w rozszerzenie elementarne $\mathfrak A_1$ struktury
$\mathfrak A_0$. Jak w dowodzie kryterium Robinsona
\ref{thm:mod-mc-robinson} utożsamiamy obraz
$\mathfrak B_0$ z nią samą, zachowując inkluzję
elementarną $\mathfrak A_0\preccurlyeq\mathfrak A_1$.
Ponieważ $\mathfrak A_1$ nadal spełnia
$T_{\forall\exists}$, możemy powtórzyć konstrukcję:
\[
\mathfrak A_0\subseteq\mathfrak B_0
\subseteq\mathfrak A_1\subseteq\mathfrak B_1
\subseteq\mathfrak A_2\subseteq\cdots,
\qquad
\mathfrak A_i\preccurlyeq\mathfrak A_{i+1},
\quad \mathfrak B_i\models T.
\]
Aby formalnie uzasadnić rekurencyjne wybory,
stosujemy dolne twierdzenie Löwenheima--Skolema
do modeli otrzymanych z obu lematów, zachowując
uniwersum poprzedniej struktury. Można ograniczyć
moce wszystkich członów przez
$\max(\aleph_0,|\mathcal L|,|A_0|)$ na mocy
twierdzenia \ref{thm:downward-ls}.
Podstruktura elementarna zawierająca poprzedni
człon zachowuje jego domknięcie egzystencjalne,
gdy zachodziło ono w większym modelu, ponieważ
formuły egzystencjalne z jego parametrami mają
w obu rozszerzeniach tę samą prawdziwość.
Przenosimy nowe elementy do ustalonych rozłącznych
zbiorów odpowiedniej mocy. Jak w dowodzie kryterium
Robinsona wybory odbywają się wtedy w obrębie zbioru.

Niech $\mathfrak U$ będzie sumą przeplatających
się łańcuchów. Z założenia (iii), zastosowanego
do ciągu $(\mathfrak B_i)$, mamy $\mathfrak U\models T$.
Z lematu \ref{lem:mod-saturation-chain}, zastosowanego
do łańcucha elementarnego $(\mathfrak A_i)$,
mamy $\mathfrak A_0\preccurlyeq\mathfrak U$.
Stąd $\mathfrak A_0\models T$. Dowolny model
$T_{\forall\exists}$ spełnia więc $T$, a odwrotna
inkluzja klas modeli wynika z definicji konsekwencji.
Otrzymujemy (i).
\end{proof}

\begin{wniosek}[Zupełność modelowa a aksjomatyzacja]
\label{cor:mod-mc-inductive}
Każda spełnialna teoria modelowo zupełna ma
aksjomatyzację uniwersalno-egzystencjalną.
\end{wniosek}
\begin{proof}
W rosnącym ciągu modeli teorii modelowo zupełnej
każda inkluzja jest elementarna. Z lematu
\ref{lem:mod-saturation-chain} suma ciągu jest
rozszerzeniem elementarnym pierwszego modelu,
więc także modelem teorii. Stosujemy twierdzenie
\ref{thm:mod-mc-inductive}.
\end{proof}

\begin{przyklad}[Implikacja odwrotna jest fałszywa]
\label{ex:mod-mc-inductive-not-mc}
Teoria wszystkich porządków liniowych w języku $\{<\}$
jest uniwersalna: aksjomaty irrefleksyjności,
przechodniości i porównywalności różnych elementów
mają postać uniwersalną z macierzami bez kwantyfikatorów.
Jest więc również uniwersalno-egzystencjalna,
z pustym blokiem egzystencjalnym.
Nie jest modelowo zupełna: inkluzja
$(\mathbb Z,<)\subseteq(\Q,<)$ nie jest elementarna,
gdyż w $\Q$ istnieje element między $0$ i $1$,
a w $\mathbb Z$ takiego elementu nie ma.
\end{przyklad}

\begin{przyklad}[Teoria uniwersalno-egzystencjalna nieuniwersalna]
\label{ex:mod-mc-inductive-not-universal}
Teoria gęstych porządków liniowych bez elementów
krańcowych jest uniwersalno-egzystencjalna. Oprócz
uniwersalnych aksjomatów porządku liniowego wystarczą
zdania
\[
\begin{gathered}
\forall x\,\forall y\,\exists z\,
(\neg(x<y)\lor(x<z\land z<y)),\\
\forall x\,\exists u\,\exists v\,(u<x\land x<v).
\end{gathered}
\]
Pierwsze wyraża gęstość, ponieważ jeśli $x<y$,
wymaga świadka między nimi, a w przeciwnym razie
dowolny element jest świadkiem. Drugie wyraża brak
elementów krańcowych. Teoria nie jest uniwersalna:
jednopunktowa podstruktura $(\Q,<)$ nie spełnia
aksjomatu braku elementów krańcowych.
Z twierdzenia \ref{thm:mod-qe-dlo} oraz wniosku
\ref{cor:mod-mc-qe} teoria ta jest modelowo zupełna.
\end{przyklad}

\section{Istnienie rozszerzeń egzystencjalnie domkniętych}

\begin{twierdzenie}[Istnienie modelu egzystencjalnie domkniętego]
\label{thm:mod-mc-ec-extension}
Niech $T$ będzie spełnialną teorią uniwersalno-egzystencjalną
w języku $\mathcal L$, a $\mathfrak M\models T$.
Istnieje rozszerzenie $\mathfrak E\supseteq\mathfrak M$,
które jest egzystencjalnie domkniętym modelem $T$.
Można przy tym zapewnić ograniczenie
\[
|E|\leq\max(\aleph_0,|\mathcal L|,|M|).
\]
Nie twierdzimy, że inkluzja
$\mathfrak M\subseteq\mathfrak E$ jest elementarna
ani że jest egzystencjalnie domknięta.
\end{twierdzenie}
\begin{proof}
Niech $\mu=\max(\aleph_0,|\mathcal L|,|M|)$.
Najpierw opiszemy pojedynczy etap dla dowolnego modelu
$\mathfrak A\models T$ mocy co najwyżej $\mu$.
Rozpatrujemy wszystkie zdania egzystencjalne w języku
$\mathcal L(A)$, otrzymanym przez dodanie nazw elementów
$A$. Obejmują one każdą formułę egzystencjalną języka
$\mathcal L$ po zastąpieniu jej zmiennych wolnych
nazwami dowolnej krotki z $A$. Jest ich co najwyżej $\mu$:
rozszerzony alfabet ma moc co najwyżej $\mu$, a każda
formuła jest skończonym ciągiem symboli, więc stosujemy
lemat \ref{lem:ls-formula-count} do języka z nazwami.
Ustalamy wyliczenie $(\varphi_\alpha)_{\alpha<\tau}$
tych formuł, gdzie $\tau\leq\mu$ jest liczbą porządkową
(utożsamiamy moc $\mu$ z jej początkową liczbą porządkową).
Przy ocenie zdań w rozszerzeniach nazwy zawsze
oznaczają pierwotne elementy $A$.

Konstruujemy rosnący łańcuch
$(\mathfrak A_\alpha)_{\alpha\leq\tau}$ modeli $T$
o mocy co najwyżej $\mu$, zaczynając od
$\mathfrak A_0=\mathfrak A$. Na etapie następnikowym
rozpatrujemy $\varphi_\alpha$.
Jeżeli istnieje rozszerzenie $\mathfrak B$ modelu
$\mathfrak A_\alpha$, będące modelem $T$, w którym
$\varphi_\alpha$ jest prawdziwa, wybieramy takie
rozszerzenie jako $\mathfrak A_{\alpha+1}$,
z dodatkowym ograniczeniem mocy.
Można je zapewnić: na mocy dolnego twierdzenia
Löwenheima--Skolema \ref{thm:downward-ls} wybieramy
podstrukturę elementarną $\mathfrak B$ zawierającą
$A_\alpha$ i o mocy co najwyżej $\mu$.
Ponieważ jest elementarna, spełnia $T$ oraz
$\varphi_\alpha$ przy nazwanych parametrach.
Jeżeli żadne rozszerzenie będące modelem $T$
nie spełnia $\varphi_\alpha$, kładziemy
$\mathfrak A_{\alpha+1}=\mathfrak A_\alpha$.

Na niezerowym etapie granicznym $\beta\leq\tau$
przyjmujemy
\[
\mathfrak A_\beta=\bigcup_{\alpha<\beta}\mathfrak A_\alpha.
\]
Jest to model $T$ na mocy twierdzenia
\ref{thm:mod-mc-inductive}. Jego moc nie przekracza
$\mu$, ponieważ $\beta\leq\mu$ oraz
$\mu\cdot\mu=\mu$ dla nieskończonej mocy $\mu$.

Wyjaśnijmy również, że wybory w tej rekurencji
można ograniczyć do zbioru. Przydzielamy każdemu
etapowi następnikowemu rozłączny zbiór nowych
elementów mocy $\mu$, rozłączny z $A$.
Jeśli wybieramy właściwe rozszerzenie, przenosimy
je na izomorficzną kopię ustalającą $A_\alpha$
i umieszczającą pozostałe elementy w przydzielonym
zbiorze. Jest to możliwe dzięki ograniczeniu mocy.
Cała konstrukcja zachodzi w obrębie struktur
na podzbiorach ustalonej sumy zbiorów, która ma
moc co najwyżej $\mu$. Struktury danego języka
na podzbiorach ustalonego zbioru tworzą zbiór:
interpretacje symboli są funkcjami i relacjami
na odpowiednich skończonych potęgach tego zbioru.
Możemy więc dobrze uporządkować zbiór kandydatów
i na każdym etapie wybierać pierwszy dopuszczalny
model, korzystając z aksjomatu wyboru i rekurencji
pozaskończonej. Wyliczenie formuł także odbywa się
w zbiorze. Nie dokonujemy wyboru z klasy właściwej.

Otrzymany model $\mathfrak A'=\mathfrak A_\tau$
ma następującą własność: jeżeli formuła egzystencjalna
$\varphi$ z parametrami z $A$ jest prawdziwa
w jakimś rozszerzeniu $\mathfrak A'$ będącym modelem $T$,
to jest prawdziwa już w $\mathfrak A'$.
Istotnie, niech $\alpha$ będzie indeksem tej formuły.
Wskazane rozszerzenie $\mathfrak A'$ jest również
rozszerzeniem $\mathfrak A_\alpha$, więc na etapie
$\alpha$ istniał dopuszczalny model realizujący
$\varphi_\alpha$. Został wybrany model
$\mathfrak A_{\alpha+1}$, w którym formuła była
prawdziwa. Ponieważ jest egzystencjalna, jej
prawdziwość zachowuje się aż do $\mathfrak A'$
na mocy lematu \ref{lem:mod-mc-preservation}.
To dowodzi własności pojedynczego etapu.

Jednokrotne wykonanie etapu nie wystarcza, ponieważ
w $\mathfrak A'$ mogą pojawić się nowe parametry.
Powtarzamy zatem etap dla kolejnych modeli, poczynając
od $\mathfrak M^{(0)}=\mathfrak M$:
\[
\mathfrak M^{(0)}\subseteq\mathfrak M^{(1)}
\subseteq\mathfrak M^{(2)}\subseteq\cdots,
\]
gdzie $\mathfrak M^{(n+1)}$ otrzymujemy z
$\mathfrak M^{(n)}$ przez opisany etap. W każdym etapie
stosujemy osobny zapas świeżych elementów; łącznie
wystarcza zbiór mocy $\aleph_0\cdot\mu\cdot\mu=\mu$.
Dobrze porządkując zbiór kandydatów oraz zbiór
formuł z nazwami z całego zapasu i $M$, uzasadniamy
wszystkie wybory w tej konstrukcji przeliczalnej.
Kładziemy
\[
\mathfrak E=\bigcup_{n<\omega}\mathfrak M^{(n)}.
\]
Ponownie z twierdzenia \ref{thm:mod-mc-inductive}
wynika $\mathfrak E\models T$ i mamy $|E|\leq\mu$.

Niech $\mathfrak N\models T$ będzie dowolnym
rozszerzeniem $\mathfrak E$, a $\varphi(\bar a)$
formułą egzystencjalną z parametrami z $E$,
prawdziwą w $\mathfrak N$.
Skończona krotka $\bar a$ leży w pewnym
$M^{(n)}$. Ponieważ $\mathfrak N$ rozszerza również
$\mathfrak M^{(n+1)}$, własność pojedynczego etapu
zapewnia prawdziwość $\varphi(\bar a)$ już w
$\mathfrak M^{(n+1)}$. Lemat
\ref{lem:mod-mc-preservation} przenosi ją do
$\mathfrak E$. Z definicji
\ref{def:mod-mc-existential-closure} model
$\mathfrak E$ jest egzystencjalnie domknięty w
każdym swoim rozszerzeniu będącym modelem $T$.
\end{proof}

\begin{figure}[htbp]
\centering
\begin{tikzpicture}[>=Stealth,x=1cm,y=1cm]
\foreach \x/\i in {0/0,3/1,6/2}{
 \node[draw,rounded corners,minimum width=1.5cm]
  (m\i) at (\x,0) {$\mathfrak M^{(\i)}$};
}
\draw[->,thick] (m0)--node[above,font=\small] {etap 0}(m1);
\draw[->,thick] (m1)--node[above,font=\small] {etap 1}(m2);
\node at (7.6,0) {$\cdots$};
\node[draw,rounded corners] (e) at (9,0) {$\mathfrak E$};
\draw[->,thick,dashed] (m2)--(7.1,0);
\draw[->,thick,dashed] (8,0)--(e);
\node[align=center,text width=3cm,font=\small] at (1.5,-1)
 {formuły z parametrami z $M^{(0)}$};
\node[align=center,text width=3cm,font=\small] at (4.5,-1)
 {formuły z parametrami z $M^{(1)}$};
\end{tikzpicture}
\caption{Kolejne etapy konstrukcji rozszerzenia
egzystencjalnie domkniętego. W każdym etapie obsługujemy
formuły z parametrami z modelu początkowego tego etapu.}
\label{fig:mod-mc-ec-construction}
\end{figure}

Rysunek \ref{fig:mod-mc-ec-construction} ilustruje
zewnętrzny ciąg konstrukcji z twierdzenia
\ref{thm:mod-mc-ec-extension}. Każda formuła egzystencjalna
z parametrami z modelu sumy jest rozpatrywana na pewnym
skończonym etapie zewnętrznym.

\begin{uwaga}[Dlaczego nie realizujemy wszystkich formuł naraz?]
\label{rem:mod-mc-ec-compatible}
Formuła może być realizowalna w pewnym rozszerzeniu
danego modelu, ale nie w każdym dalszym rozszerzeniu
innego modelu uzyskanego po wcześniejszych wyborach.
Nie zakładamy wspólnego rozszerzenia wszystkich
kandydatów. Dla przykładu rozważmy teorię uniwersalną
w języku z relacją dwuargumentową $R$ i aksjomatem
\[
\forall u\,\forall v\,\forall w\,\forall z\,
((R(u,v)\land R(w,z))\to v=z).
\]
Niech model początkowy ma dwa różne elementy $a,b$
i pustą relację $R$. Formuły
$\exists u\,R(u,a)$ i $\exists u\,R(u,b)$
można zrealizować osobno, dodając świeży element
$u$ i odpowiednią pojedynczą parę do relacji.
Nie można ich jednak zrealizować łącznie,
bo aksjomat wymuszałby $a=b$.
Na każdym etapie pytamy więc o realizowalność
w rozszerzeniu \emph{aktualnego} modelu.
Wystarcza to, ponieważ rozszerzenie modelu końcowego
jest także rozszerzeniem wszystkich modeli pośrednich.
\end{uwaga}

\begin{przyklad}[Bez hipotezy o sumach łańcuchów]
\label{ex:mod-mc-no-ec}
Niech $T$ będzie teorią porządków liniowych
mających element największy, w języku $\{<\}$.
Teoria jest spełnialna, lecz nie ma żadnego
egzystencjalnie domkniętego modelu.
Jeśli $\mathfrak A\models T$ i $a$ jest jego
elementem największym, dodajemy świeży element
$t$ większy od wszystkich elementów $A$.
Otrzymujemy model $T$ rozszerzający $\mathfrak A$,
w którym $\exists y\,(a<y)$ jest prawdziwa,
chociaż jest fałszywa w $\mathfrak A$.
To wyklucza domknięcie egzystencjalne dowolnego
modelu $T$.

Nie przeczy to twierdzeniu
\ref{thm:mod-mc-ec-extension}, gdyż teoria nie jest
uniwersalno-egzystencjalna. Istotnie, porządki
na zbiorach $\{0,\ldots,n\}$ dla $n\in\N$
są modelami $T$ i tworzą rosnący ciąg podstruktur,
ale ich suma $(\N,<)$ nie ma elementu największego.
Stosujemy twierdzenie \ref{thm:mod-mc-inductive}.
\end{przyklad}

\begin{przyklad}[Domknięcie egzystencjalne porządków liniowych]
\label{ex:mod-mc-ec-dlo}
Egzystencjalnie domknięte modele teorii wszystkich
porządków liniowych są dokładnie gęstymi porządkami
liniowymi bez elementów krańcowych.
\begin{proof}
Niech $\mathfrak A$ będzie porządkiem liniowym.
Możemy zanurzyć go w gęsty porządek bez elementów
krańcowych na $A\times\Q$ z porządkiem leksykograficznym:
\[
(a,q)<(b,r)
\quad\Longleftrightarrow\quad
a<b\ \lor\ (a=b\land q<r).
\]
Jest to porządek liniowy, ponieważ o porównaniu
decyduje pierwsza różna współrzędna, a liniowość
i przechodniość porządków obu współrzędnych
zapewniają te własności porządku par.
Jeżeli pierwsze współrzędne są równe, między
dwoma różnymi punktami wybieramy punkt o tej samej
pierwszej współrzędnej i wymiernej drugiej współrzędnej
pomiędzy danymi. Jeżeli $a<b$, to między
$(a,q)$ i $(b,r)$ leży $(a,q+1)$.
Brak elementów krańcowych wynika z istnienia
$(a,q-1)$ i $(a,q+1)$ dla każdego $(a,q)$.
Odwzorowanie $a\mapsto(a,0)$ jest zanurzeniem;
przenosimy strukturę par na rzeczywiste rozszerzenie
$\mathfrak A$.

Jeżeli $\mathfrak A$ jest egzystencjalnie domknięta
w klasie porządków liniowych, z powyższego rozszerzenia
przenosimy formuły
$\exists z\,(a<z\land z<b)$ dla $a<b$ oraz
$\exists u\,(u<a)$ i $\exists v\,(a<v)$.
Dają one gęstość i brak elementów krańcowych.

Odwrotnie, niech $\mathfrak A$ będzie gęstym porządkiem
bez elementów krańcowych i niech $\mathfrak B$
będzie dowolnym rozszerzającym ją porządkiem liniowym.
Opisaną konstrukcją rozszerzamy $\mathfrak B$
do gęstego porządku bez elementów krańcowych
$\mathfrak C$. Obie struktury $\mathfrak A$ i
$\mathfrak C$ są modelami tej samej teorii z eliminacją
kwantyfikatorów (twierdzenie \ref{thm:mod-qe-dlo}),
więc $\mathfrak A\preccurlyeq\mathfrak C$ na mocy
wniosku \ref{cor:mod-mc-qe}. Jeżeli formuła egzystencjalna
z parametrami z $A$ zachodzi w $\mathfrak B$, zachodzi
w $\mathfrak C$ przez lemat \ref{lem:mod-mc-preservation},
a następnie w $\mathfrak A$ przez elementarność.
To daje żądane domknięcie egzystencjalne.
\end{proof}
\end{przyklad}

\begin{przyklad}[Grafy i własność przedłużania]
\label{ex:mod-mc-ec-graphs}
Niech $T_{\mathrm G}$ będzie teorią grafów prostych
nieskierowanych w języku z dwuargumentową relacją $R$,
gdzie aksjomaty wymagają symetrii i irrefleksyjności.
Model $\mathfrak G\models T_{\mathrm G}$ jest
egzystencjalnie domknięty wtedy i tylko wtedy,
gdy dla dowolnych rozłącznych skończonych podzbiorów
$U,V\subseteq G$ istnieje $w\in G\setminus(U\cup V)$,
który jest połączony z każdym elementem $U$
i nie jest połączony z żadnym elementem $V$.
\begin{proof}
Aksjomaty symetrii i irrefleksyjności są uniwersalne.
Jeśli $\mathfrak G$ jest egzystencjalnie domknięty,
możemy rozpatrzyć rozszerzenie grafu przez dodanie
jednego świeżego wierzchołka połączonego dokładnie
z wierzchołkami $U$ (i z żadnymi innymi starymi).
Relację rozszerzamy symetrycznie, nie dodając pętli.
W rozszerzeniu zachodzi formuła egzystencjalna
\[
\exists w\left(
\bigwedge_{a\in U\cup V}w\ne a
\land\bigwedge_{a\in U}R(w,a)
\land\bigwedge_{b\in V}\neg R(w,b)
\right).
\]
Zatem zachodzi także w $\mathfrak G$.

Odwrotnie, załóżmy własność przedłużania i niech
$\mathfrak H\models T_{\mathrm G}$ rozszerza $\mathfrak G$.
Niech $\exists\bar y\,\theta(\bar a,\bar y)$,
gdzie $\theta$ nie zawiera kwantyfikatorów i
$\bar a$ pochodzi z $G$, będzie prawdziwa w $\mathfrak H$.
Wybieramy krotkę świadków $\bar b$ w $\mathfrak H$.
Chcemy przenieść skończony graf indukowany przez
współrzędne $(\bar a,\bar b)$ do $\mathfrak G$,
ustalając współrzędne $\bar a$.

Najpierw ustalamy identyczności wszystkich różnych
elementów występujących w $\bar a$.
Pozostałe różne elementy krotki $\bar b$ rozpatrujemy
po kolei. Załóżmy, że wybrano już różne obrazy
dotychczasowych elementów i zachowano wszystkie
krawędzie oraz ich brak. Dla kolejnego elementu $b$
dzielimy dotychczasowe obrazy na zbiór $U$ obrazów
sąsiadów $b$ i zbiór $V$ obrazów pozostałych elementów.
Własność przedłużania daje nowy wierzchołek różny
od wszystkich dotychczasowych, o dokładnie wymaganej
relacji do nich. W ten sposób rozszerzamy zanurzenie
skończonego grafu. Powtarzające się współrzędne
wysyłamy na ten sam obraz.
Po skończonej liczbie kroków otrzymujemy krotkę
$\bar b'$ w $G$ o tym samym układzie równości
i relacji $R$ co $\bar b$ razem z $\bar a$.
Ponieważ w języku grafów są to wszystkie atomy,
a stałe logiczne mają ustalone wartości, każda formuła bez
kwantyfikatorów zachowuje wartość logiczną.
W szczególności $\mathfrak G\models\theta(\bar a,\bar b')$.
To dowodzi egzystencjalnego domknięcia.
\end{proof}
\end{przyklad}

Własność przedłużania z przykładu
\ref{ex:mod-mc-ec-graphs} można aksjomatyzować:
dla każdej pary nieujemnych liczb całkowitych
$m,n$ zapisujemy zdanie wymagające świadka dla
$m+n$ parami różnych parametrów, z których pierwsze
$m$ ma być połączonych ze świadkiem, a pozostałe
$n$ nie. W przypadku $m=n=0$ żądamy jedynie
istnienia elementu. Z twierdzenia
\ref{thm:mod-mc-ec-extension} wynika istnienie
takich grafów i możliwość rozszerzenia dowolnego
grafu prostego do grafu spełniającego schemat.
\section{Towarzysze modelowi}
\begin{definicja}[Towarzysz modelowy]
\label{def:mod-mc-companion}
Niech $T$ będzie spełnialną teorią języka $\mathcal L$.
Teorię $T^*$ tego samego języka nazywamy towarzyszem
modelowym (ang. \emph{model companion}) teorii $T$,
jeżeli $T^*$ jest spełnialna, modelowo zupełna oraz
\[
T^*_{\forall}=T_{\forall}.
\]
Równość dotyczy zbiorów konsekwencji semantycznych,
a nie wybranych zbiorów aksjomatów.
\end{definicja}

\begin{lemat}[Wzajemne zanurzenia]
\label{lem:mod-mc-mutual-embedding}
Dla dwóch spełnialnych teorii $T,S$ tego samego języka
równość $T_{\forall}=S_{\forall}$ jest równoważna
łącznemu spełnieniu następujących dwóch warunków:
\begin{enumerate}
\item każdy model $T$ zanurza się w pewien model $S$;
\item każdy model $S$ zanurza się w pewien model $T$.
\end{enumerate}
\end{lemat}
\begin{proof}
Jeżeli konsekwencje uniwersalne są identyczne,
każdy model jednej teorii spełnia uniwersalne
konsekwencje drugiej. Twierdzenie
\ref{thm:mod-mc-universal-embedding} daje oba
wymagane zanurzenia.
Odwrotnie, niech $\sigma$ będzie uniwersalną
konsekwencją $S$ i niech $\mathfrak A\models T$.
Zanurzamy $\mathfrak A$ w model $S$.
Z zachowania zdań uniwersalnych przy podstrukturach
(twierdzenie \ref{thm:model-universal-substructure})
oraz niezmienniczości względem izomorfizmów wynika
$\mathfrak A\models\sigma$. Stąd
$S_{\forall}\subseteq T_{\forall}$. Zamiana ról
teorii daje odwrotną inkluzję.
\end{proof}

Towarzysz modelowy nie musi być rozszerzeniem teorii
wyjściowej nawet w sensie semantycznym, tj. nie musi
mieć wyłącznie modeli teorii wyjściowej. Własność
wzajemnych zanurzeń mówi jedynie, że modele obu
teorii można rozszerzać do modeli drugiej teorii.
Dla teorii uniwersalno-egzystencjalnych uzyskamy
silniejszy wynik.

\begin{twierdzenie}[Jednoznaczność towarzysza modelowego]
\label{thm:mod-mc-companion-unique}
Jeśli teoria $T$ ma dwóch towarzyszy modelowych
$S_1,S_2$, to $S_1$ i $S_2$ mają te same modele.
W szczególności ich domknięcia na konsekwencję
semantyczną są identyczne.
\end{twierdzenie}
\begin{proof}
Mamy $(S_1)_{\forall}=T_{\forall}=(S_2)_{\forall}$.
Niech $\mathfrak A_0\models S_1$. Na mocy lematu
\ref{lem:mod-mc-mutual-embedding} możemy kolejno
rozszerzać modele obu teorii, uzyskując
\[
\mathfrak A_0\subseteq\mathfrak B_0
\subseteq\mathfrak A_1\subseteq\mathfrak B_1
\subseteq\cdots,
\]
gdzie $\mathfrak A_i\models S_1$, a
$\mathfrak B_i\models S_2$. W każdym kroku zamieniamy
zanurzenie na inkluzję przez przeniesienie struktury
na kopię ustalającą poprzedni model.
Dolne twierdzenie Löwenheima--Skolema
\ref{thm:downward-ls} pozwala ograniczyć moce wszystkich
członów przez $\mu=\max(\aleph_0,|\mathcal L|,|A_0|)$,
wybierając podstrukturę elementarną kolejnego modelu
zawierającą poprzedni. Wybory można więc przeprowadzić
w obrębie zbioru struktur, z rozłącznymi zapasami nowych
elementów, jak w dowodzie twierdzenia
\ref{thm:mod-mc-ec-extension}.

Z zupełności modelowej $S_1$ i $S_2$ wynika
$\mathfrak A_i\preccurlyeq\mathfrak A_{i+1}$
oraz $\mathfrak B_i\preccurlyeq\mathfrak B_{i+1}$.
Oba łańcuchy mają tę samą strukturę sumy $\mathfrak U$.
Na mocy lematu \ref{lem:mod-saturation-chain}
$\mathfrak U\models S_2$ oraz
$\mathfrak A_0\preccurlyeq\mathfrak U$.
W szczególności $\mathfrak A_0\models S_2$.
Zamiana ról $S_1,S_2$ dowodzi odwrotnej inkluzji
klas modeli.
\end{proof}

\begin{lemat}[Towarzysz teorii uniwersalno-egzystencjalnej]
\label{lem:mod-mc-companion-extension}
Jeśli $T$ jest uniwersalno-egzystencjalna i ma
towarzysza modelowego $T^*$, to każdy model $T^*$
jest modelem $T$.
\end{lemat}
\begin{proof}
Niech $\mathfrak A_0\models T^*$. Jak w poprzednim
dowodzie, wzajemne zanurzenia dają przeplatające się
łańcuchy
\[
\mathfrak A_0\subseteq\mathfrak B_0
\subseteq\mathfrak A_1\subseteq\mathfrak B_1
\subseteq\cdots
\]
z $\mathfrak A_i\models T^*$ i $\mathfrak B_i\models T$.
Moce i wybory kontrolujemy tak samo jak w dowodzie
twierdzenia \ref{thm:mod-mc-companion-unique}.
Zupełność modelowa $T^*$ zapewnia elementarność
łańcucha $(\mathfrak A_i)$.
Jego suma $\mathfrak U$ jest także sumą rosnącego
łańcucha $(\mathfrak B_i)$, którego inkluzje nie
muszą być elementarne.
Z twierdzenia \ref{thm:mod-mc-inductive} mamy jednak
$\mathfrak U\models T$. Ponieważ
$\mathfrak A_0\preccurlyeq\mathfrak U$ na mocy
lematu o sumach łańcuchów elementarnych
\ref{lem:mod-saturation-chain}, otrzymujemy
$\mathfrak A_0\models T$.
\end{proof}

\begin{twierdzenie}[Towarzysz modelowy a domknięcie egzystencjalne]
\label{thm:mod-mc-companion-ec}
Niech $T$ będzie spełnialną teorią
uniwersalno-egzystencjalną. Następujące warunki
są równoważne:
\begin{enumerate}[label=(\roman*)]
\item $T$ ma towarzysza modelowego;
\item klasa egzystencjalnie domkniętych modeli $T$
jest klasą modeli pewnej teorii pierwszego rzędu
w języku $T$.
\end{enumerate}
Jeżeli te warunki zachodzą, modele towarzysza
modelowego są dokładnie egzystencjalnie domkniętymi
modelami $T$.
\end{twierdzenie}
\begin{proof}
Załóżmy (i) i niech $T^*$ będzie towarzyszem.
Z lematu \ref{lem:mod-mc-companion-extension}
każdy model $T^*$ jest modelem $T$.

Niech $\mathfrak A\models T^*$, a
$\mathfrak B\models T$ będzie rozszerzeniem
$\mathfrak A$. Z lematu \ref{lem:mod-mc-mutual-embedding}
istnieje rozszerzenie $\mathfrak C\supseteq\mathfrak B$
będące modelem $T^*$.
Ponieważ $\mathfrak A,\mathfrak C\models T^*$,
zupełność modelowa daje
$\mathfrak A\preccurlyeq\mathfrak C$.
Jeżeli formuła egzystencjalna z parametrami z $A$
jest prawdziwa w $\mathfrak B$, jest prawdziwa
w $\mathfrak C$ na mocy lematu
\ref{lem:mod-mc-preservation}, a następnie
w $\mathfrak A$ przez elementarność.
Zatem $\mathfrak A$ jest egzystencjalnie domkniętym
modelem $T$.

Odwrotnie, niech $\mathfrak A$ będzie egzystencjalnie
domkniętym modelem $T$. Rozszerzamy ją do modelu
$\mathfrak C\models T^*$ przy użyciu wzajemnych
zanurzeń. Z lematu \ref{lem:mod-mc-companion-extension}
$\mathfrak C\models T$, więc
$\mathfrak A\preccurlyeq_{\exists}\mathfrak C$.
Teoria $T^*$ ma aksjomatyzację
uniwersalno-egzystencjalną na mocy wniosku
\ref{cor:mod-mc-inductive}.
Każdy taki aksjomat
$\forall\bar x\,\exists\bar y\,\theta(\bar x,\bar y)$
zachodzi w $\mathfrak A$: dla dowolnej krotki
$\bar a$ z $A$ formuła egzystencjalna
$\exists\bar y\,\theta(\bar a,\bar y)$ jest
prawdziwa w $\mathfrak C$, więc jest prawdziwa
w $\mathfrak A$ przez domknięcie egzystencjalne.
Dotyczy to również pustych bloków kwantyfikatorów.
Stąd $\mathfrak A\models T^*$.
Udowodniliśmy charakterystykę modeli towarzysza
i w szczególności (ii).

Załóżmy teraz (ii). Niech $S$ będzie teorią
aksjomatyzującą dokładnie egzystencjalnie domknięte
modele $T$. Jest spełnialna na mocy twierdzenia
\ref{thm:mod-mc-ec-extension}, ponieważ $T$ jest
spełnialna i uniwersalno-egzystencjalna.
Każdy model $T$ ma rozszerzenie będące modelem $S$,
a każdy model $S$ jest modelem $T$.
Lemat \ref{lem:mod-mc-mutual-embedding} daje więc
$S_{\forall}=T_{\forall}$.
Jeśli $\mathfrak A\subseteq\mathfrak B$ są modelami $S$,
to $\mathfrak A$ jest egzystencjalnie domkniętym
modelem $T$, a $\mathfrak B\models T$.
W konsekwencji $\mathfrak A\preccurlyeq_{\exists}
\mathfrak B$. Kryterium Robinsona
\ref{thm:mod-mc-robinson} daje zupełność modelową $S$.
Zatem $S$ jest towarzyszem modelowym $T$.
\end{proof}

\begin{wniosek}[Redukcja do konsekwencji uniwersalnych]
\label{cor:mod-mc-companion-universal}
Spełnialna teoria $T$ ma towarzysza modelowego dokładnie
wtedy, gdy $T_{\forall}$ ma towarzysza modelowego.
W takim przypadku są to te same teorie z dokładnością
do równoważności semantycznej.
\end{wniosek}
\begin{proof}
Każde zdanie uniwersalne wynikające z $T_{\forall}$
wynika z $T$, ponieważ każdy model $T$ spełnia
$T_{\forall}$. Odwrotnie każde zdanie z $T_{\forall}$
jest oczywiście konsekwencją $T_{\forall}$.
Stąd $(T_{\forall})_{\forall}=T_{\forall}$.
Definicja \ref{def:mod-mc-companion} daje zatem
dokładnie te same warunki dla towarzysza obu teorii,
a twierdzenie \ref{thm:mod-mc-companion-unique}
zapewnia jednoznaczność.
\end{proof}

Wniosek ten pozwala w przypadku dowolnej spełnialnej
teorii $T$ badać istnienie towarzysza przez klasę
egzystencjalnie domkniętych modeli $T_{\forall}$.
Teoria $T_{\forall}$ jest uniwersalna, więc także
uniwersalno-egzystencjalna, a zatem można zastosować
twierdzenie \ref{thm:mod-mc-companion-ec}.
Nie należy natomiast bez dodatkowych hipotez
utożsamiać tej klasy z egzystencjalnie domkniętymi
modelami samej teorii $T$.

\begin{przyklad}[Towarzysz, który nie rozszerza teorii wyjściowej]
\label{ex:mod-mc-companion-not-extension}
Niech $T$ będzie teorią porządków liniowych mających
element największy, a $S$ teorią wszystkich porządków
liniowych. Każdy model $T$ jest modelem $S$.
Każdy model $S$ można rozszerzyć do modelu $T$
przez dodanie nowego elementu większego od wszystkich
dotychczasowych. Z lematu
\ref{lem:mod-mc-mutual-embedding} wynika
$T_{\forall}=S_{\forall}$. Ponieważ $S$ jest
uniwersalna (przykład \ref{ex:mod-mc-inductive-not-mc}),
ma te same modele co $T_{\forall}$.
Egzystencjalnie domknięte modele $S$ to dokładnie
gęste porządki bez elementów krańcowych na mocy
przykładu \ref{ex:mod-mc-ec-dlo}. Z twierdzenia
\ref{thm:mod-mc-companion-ec} i wniosku
\ref{cor:mod-mc-companion-universal} ich teoria
jest towarzyszem modelowym zarówno $S$, jak i $T$.

Żaden model tego towarzysza nie jest modelem $T$,
ponieważ nie ma elementu największego.
Co więcej, jak wykazano w przykładzie
\ref{ex:mod-mc-no-ec}, sama teoria $T$ nie ma
żadnych modeli egzystencjalnie domkniętych.
Pokazuje to konieczność hipotezy o aksjomatyzacji
uniwersalno-egzystencjalnej w twierdzeniu
\ref{thm:mod-mc-companion-ec}.
\end{przyklad}

\begin{wniosek}[Zupełność towarzysza modelowego]
\label{cor:mod-mc-companion-complete}
Niech $T$ ma towarzysza modelowego $T^*$. Wówczas
następujące warunki są równoważne:
\begin{enumerate}[label=(\roman*)]
\item $T^*$ jest zupełna;
\item dowolne dwa modele $T$ można zanurzyć w jeden
model $T$;
\item dowolne dwa modele $T^*$ można zanurzyć w jeden
model $T^*$.
\end{enumerate}
Warunek (ii) nazywa się własnością wspólnego zanurzenia
(ang. \emph{joint embedding property}).
\end{wniosek}
\begin{proof}
Załóżmy (i). Dwa modele $T^*$ zgadzają się co do
wszystkich zdań, w szczególności egzystencjalnych.
Lemat \ref{lem:mod-mc-existential-type} dla pustych
krotek daje zanurzenie pierwszego modelu w rozszerzenie
elementarne drugiego. To rozszerzenie jest modelem
$T^*$ i zawiera drugi model, więc zachodzi (iii).

Załóżmy (iii). Dwa dowolne modele $T$ zanurzamy
najpierw w modele $T^*$, korzystając z lematu
\ref{lem:mod-mc-mutual-embedding}. Z (iii) te dwa
modele zanurzają się w jeden model $T^*$, który
następnie zanurzamy w model $T$. Złożenia zanurzeń
dają (ii).

Załóżmy (ii). Odwracamy role teorii w poprzednim
argumencie: dwa modele $T^*$ zanurzamy w modele $T$,
te w jeden model $T$, a ten w model $T^*$.
Otrzymujemy (iii). Ponieważ $T^*$ jest modelowo
zupełna, oba zanurzenia z (iii) są elementarne,
a zatem ich dziedziny zgadzają się co do wszystkich
zdań, których prawdziwość można odczytać we wspólnym
modelu. Wobec spełnialności $T^*$ daje to (i).
\end{proof}

\begin{przyklad}[Towarzysz teorii porządków liniowych]
\label{ex:mod-mc-companion-dlo}
Towarzyszem modelowym teorii wszystkich porządków
liniowych jest teoria gęstych porządków liniowych
bez elementów krańcowych. Istotnie, teoria porządków
liniowych jest uniwersalna, a w przykładzie
\ref{ex:mod-mc-ec-dlo} scharakteryzowaliśmy jej
egzystencjalnie domknięte modele. Twierdzenie
\ref{thm:mod-mc-companion-ec} daje tezę.
\end{przyklad}

\begin{przyklad}[Towarzysz teorii grafów prostych]
\label{ex:mod-mc-companion-graphs}
Niech $T_{\mathrm G}$ będzie teorią grafów prostych
nieskierowanych z przykładu \ref{ex:mod-mc-ec-graphs}.
Jej towarzyszem modelowym jest teoria
$T_{\mathrm G}^*$ otrzymana przez dodanie dla wszystkich
$m,n\geq0$ aksjomatów przedłużania:
\[
\begin{aligned}
\forall\bar u\,\forall\bar v\,\exists w\,\bigg(
 &\neg D(\bar u,\bar v)\\
 {}\lor{}&
 \Big(\bigwedge_{a\in(\bar u,\bar v)}w\ne a
 \land\bigwedge_{i=1}^{m}R(w,u_i)\\
 &{}\land\bigwedge_{j=1}^{n}\neg R(w,v_j)\Big)\bigg).
\end{aligned}
\]
Tutaj $\bar u=(u_1,\ldots,u_m)$,
$\bar v=(v_1,\ldots,v_n)$ i $D(\bar u,\bar v)$
oznacza koniunkcję nierówności między wszystkimi
parami różnych pozycji połączonej listy zmiennych.
Koniunkcja indeksowana przez $(\bar u,\bar v)$
obejmuje wszystkie jej pozycje, nie zaś elementy
jakiegoś ustalonego grafu.
Jeśli lista parametrów nie ma powtórzeń, aksjomat
wymaga świadka o odpowiednim układzie sąsiedztwa.
W przeciwnym razie dowolny element jest świadkiem,
ponieważ uniwersum jest niepuste.
Dla pustej listy aksjomat sprowadza się do
$\exists w\,\top$.

Modele teorii $T_{\mathrm G}$ wraz z tym schematem
są dokładnie egzystencjalnie domkniętymi grafami
na mocy przykładu
\ref{ex:mod-mc-ec-graphs}. Ponieważ $T_{\mathrm G}$
jest uniwersalna, twierdzenie
\ref{thm:mod-mc-companion-ec} daje istnienie
towarzysza modelowego $T_{\mathrm G}^*$.
\end{przyklad}

\begin{wniosek}[Eliminacja kwantyfikatorów w teorii grafu Rado]
\label{cor:mod-mc-rado-qe}
Teoria $T_{\mathrm G}^*$ ma eliminację kwantyfikatorów
i jest zupełna. Ma też model przeliczalny.
\end{wniosek}
\begin{proof}
Niech $\mathfrak A,\mathfrak B\models T_{\mathrm G}^*$
oraz niech $\bar a,\bar b$ mają ten sam typ
bez kwantyfikatorów. Krotki mają ten sam układ
równości i krawędzi, ponieważ w tym języku atomy
są równościami zmiennych i formułami $R(x_i,x_j)$.
Dla $a\in A$ należącego do krotki $\bar a$
wybieramy odpowiadającą współrzędną $\bar b$.
Wybór jest spójny dzięki zgodności równości.
Jeżeli $a$ jest nowym elementem, dzielimy zbiór
różnych współrzędnych $\bar b$ na zbiór $U$ obrazów
sąsiadów $a$ i zbiór $V$ obrazów pozostałych
współrzędnych. Aksjomat przedłużania daje
$b\in B\setminus(U\cup V)$ o wymaganym układzie
krawędzi. Symetria i irrefleksyjność relacji
zapewniają także zgodność atomów z zamienioną
kolejnością argumentów i atomów postaci $R(a,a)$.
Rozszerzone krotki mają więc ten sam typ
bez kwantyfikatorów. Gdy krotki są puste,
wybieramy dowolny element $B$.
Twierdzenie \ref{thm:mod-qe-forth} daje eliminację
kwantyfikatorów.

W języku grafów nie ma stałych, symboli funkcyjnych
ani relacji o arności zero, więc każde zdanie
bez kwantyfikatorów jest kombinacją logiczną
$\top,\bot$. Na mocy wniosku
\ref{cor:mod-qe-consequences} teoria jest zupełna.
Istnienie modelu o mocy co najwyżej $\aleph_0$
wynika z twierdzenia \ref{thm:mod-mc-ec-extension}
zastosowanego do jednopunktowego grafu.
Każdy model schematu jest nieskończony: dla dowolnego
skończonego zbioru $V$ stosujemy własność przedłużania
z $U=\varnothing$, otrzymując element poza $V$.
Uzyskany model jest zatem przeliczalnie nieskończony.

Można również wykazać zupełność bez odwoływania się
do eliminacji kwantyfikatorów: dwa dowolne grafy proste
zanurzają się w ich sumę rozłączną, w której nie dodajemy
krawędzi pomiędzy składnikami. Relacja na każdym
składniku pozostaje niezmieniona, więc są to zanurzenia.
Teoria $T_{\mathrm G}$ ma zatem własność wspólnego
zanurzenia i wniosek \ref{cor:mod-mc-companion-complete}
daje zupełność jej towarzysza.
\end{proof}

Przeliczalny model teorii $T_{\mathrm G}^*$ jest
nazywany grafem Rado. Jego jednoznaczność z dokładnością
do izomorfizmu wykażemy przy omawianiu kategoryczności.
Tutaj istotne jest to, że własność przedłużania
jednocześnie opisuje domknięcie egzystencjalne,
aksjomatyzuje towarzysza modelowego i umożliwia
eliminację kwantyfikatorów.

\section{Uzupełnienia modelowe i amalgamacja}
W tej części przyjmujemy, że $T$ jest spełnialną teorią
uniwersalną. W przypadku ogólniejszych teorii rozważa
się również inne konwencje definicji uzupełnienia
modelowego. Poniżej jawnie rozróżniamy zupełność nad
niepustym modelem $T$ od eliminacji kwantyfikatorów
obejmującej także zdania w językach bez termów domkniętych.

\begin{definicja}[Uzupełnienie modelowe]
\label{def:mod-mc-completion}
Towarzysza modelowego $T^*$ teorii uniwersalnej $T$
nazywamy jej uzupełnieniem modelowym
(ang. \emph{model completion}), jeżeli dla każdego
$\mathfrak A\models T$ teoria
\[
T^*\cup\operatorname{Diag}(\mathfrak A)
\]
jest zupełna w języku rozszerzonym o nazwy elementów $A$.
Zupełność rozumiemy wraz ze spełnialnością tej teorii.
\end{definicja}

Z lematu \ref{lem:mod-mc-companion-extension}
wynika, że każdy model $T^*$ jest modelem $T$,
ponieważ teoria uniwersalna jest także
uniwersalno-egzystencjalna.
Spełnialność teorii z diagramem jest zapewniona
przez wzajemne zanurzenia
(lemat \ref{lem:mod-mc-mutual-embedding}).
Warunek zupełności oznacza, że wszystkie rozszerzenia
danego modelu $T$ do modeli $T^*$ są elementarnie
równoważne w języku z nazwami tego modelu.
W szczególności nie wymagamy izomorfizmu takich
rozszerzeń ani ich jednakowej mocy.

\begin{definicja}[Własność amalgamacji]
\label{def:mod-mc-amalgamation}
Mówimy, że teoria $T$ ma własność amalgamacji, jeżeli
dla dowolnych modeli $\mathfrak A,\mathfrak B,\mathfrak C$
teorii $T$ i zanurzeń
$f:\mathfrak A\to\mathfrak B$, $g:\mathfrak A\to\mathfrak C$
istnieją model $\mathfrak D\models T$ i zanurzenia
$h:\mathfrak B\to\mathfrak D$, $k:\mathfrak C\to\mathfrak D$
takie, że $h\circ f=k\circ g$.
Jeżeli można dodatkowo wymagać
\[
h(B)\cap k(C)=h(f(A)),
\]
mówimy o silnej własności amalgamacji.
\end{definicja}

Dzięki zastąpieniu struktur izomorficznymi kopiami
możemy w dowodach zakładać, że $f,g$ są inkluzjami
oraz $B\cap C=A$. Robimy to, pozostawiając elementy
$A$ bez zmian, a pozostałe elementy obu uniwersów
zastępując rozłącznymi zbiorami świeżych elementów.
Zanurzenia amalgamujące muszą się zgadzać na $A$,
lecz w zwykłej amalgamacji mogą utożsamiać również
inne elementy obu rozszerzeń.

\begin{figure}[htbp]
\centering
\begin{tikzpicture}[>=Stealth]
\node (a) at (0,0) {$\mathfrak A$};
\node (b) at (-2,1.5) {$\mathfrak B$};
\node (c) at (2,1.5) {$\mathfrak C$};
\node (d) at (0,3) {$\mathfrak D$};
\draw[->,thick] (a)--node[left] {$f$}(b);
\draw[->,thick] (a)--node[right] {$g$}(c);
\draw[->,thick,blue!65!black] (b)--node[left] {$h$}(d);
\draw[->,thick,blue!65!black] (c)--node[right] {$k$}(d);
\end{tikzpicture}
\caption{Amalgamacja dwóch zanurzeń. Obie drogi
z $\mathfrak A$ do $\mathfrak D$ dają to samo zanurzenie.}
\label{fig:mod-mc-amalgamation}
\end{figure}

Rysunek \ref{fig:mod-mc-amalgamation} przedstawia
warunek przemienności z definicji
\ref{def:mod-mc-amalgamation}. Nie jest to żądanie
uniwersalnej własności koproduktu z utożsamieniem:
wymagamy istnienia wspólnego modelu i zanurzeń,
a nie jego jednoznaczności.

\begin{twierdzenie}[Uzupełnienie modelowe a amalgamacja]
\label{thm:mod-mc-completion-ap}
Niech $T$ będzie teorią uniwersalną mającą towarzysza
modelowego $T^*$. Wówczas $T^*$ jest uzupełnieniem
modelowym $T$ dokładnie wtedy, gdy $T$ ma własność
amalgamacji.
\end{twierdzenie}
\begin{proof}
Załóżmy własność amalgamacji i ustalmy
$\mathfrak A\models T$. Rozważmy dwa modele
teorii $T^*\cup\operatorname{Diag}(\mathfrak A)$.
Na mocy lematu \ref{lem:mod-mc-diagram-embedding}
ich redukty $\mathfrak B,\mathfrak C$ zawierają
zanurzone kopie $\mathfrak A$. Po przeniesieniu
struktur przyjmujemy, że oba zawierają dosłownie
$\mathfrak A$ i nazwy elementów $A$ są interpretowane
standardowo. Ponieważ $\mathfrak B,\mathfrak C$
są modelami $T$, amalgamujemy je nad $\mathfrak A$
w pewnym modelu $\mathfrak D\models T$.
Rozszerzamy $\mathfrak D$ do modelu
$\mathfrak E\models T^*$, korzystając z lematu
\ref{lem:mod-mc-mutual-embedding}.
Złożenia zanurzeń z inkluzją w $\mathfrak E$
są zanurzeniami modeli $T^*$, więc są elementarne.
Zgadzają się one na $A$.
Każde zdanie w języku z nazwami elementów $A$
ma więc tę samą wartość logiczną w obu modelach:
można ją odczytać w $\mathfrak E$ przy wspólnych
obrazach elementów $A$. Ponieważ teoria z diagramem
jest spełnialna, jest zupełna.

Odwrotnie, załóżmy zupełność teorii
$S=T^*\cup\operatorname{Diag}(\mathfrak A)$
dla każdego modelu $\mathfrak A$ teorii $T$.
Ustalmy rozszerzenia $\mathfrak B,\mathfrak C$ modelu
$\mathfrak A$, będące modelami $T$.
Nazwy elementów $A$ wybieramy wspólne w diagramach,
a nazwy elementów spoza $A$ wybieramy rozłącznie.
Twierdzimy, że teoria
\[
S\cup\operatorname{Diag}(\mathfrak B)
\cup\operatorname{Diag}(\mathfrak C)
\]
jest spełnialna. Każdy z dwóch diagramów osobno jest
zgodny z $S$, ponieważ odpowiedni model rozszerza
się do modelu $T^*$.

Niech $\theta_B(\bar c,\bar d)$ i
$\theta_C(\bar c,\bar e)$ będą koniunkcjami dowolnych
skończonych fragmentów obu diagramów. Nazwy $\bar c$
pochodzą z $A$, natomiast listy $\bar d,\bar e$
są rozłączne i nie występują w języku $S$.
Zgodność każdego diagramu z $S$ oznacza, że
w pewnych modelach $S$ zachodzą zdania
\[
\exists\bar y\,\theta_B(\bar c,\bar y),
\qquad
\exists\bar z\,\theta_C(\bar c,\bar z)
\]
odpowiednio. Ponieważ $S$ jest zupełna, oba zdania
zachodzą w każdym modelu $S$. W jednym takim modelu
wybieramy świadków obu formuł i interpretujemy nimi
nazwy $\bar d,\bar e$. Rozłączność list nazw zapewnia
niezależność tych interpretacji. Pozostałe dodatkowe
nazwy interpretujemy dowolnie. Każdy skończony fragment
sumy teorii i diagramów ma zatem model.
Zwartość daje model całej sumy, a własność diagramów
daje zanurzenia $\mathfrak B,\mathfrak C$
zgodne na $A$. Redukt wspólnego modelu spełnia $T^*$,
a więc także $T$. To jest żądana amalgamacja.
\end{proof}

\section{Amalgamacja a eliminacja kwantyfikatorów}
\begin{lemat}[Zupełność nad nazwanymi elementami]
\label{lem:mod-mc-qe-completion}
Jeżeli $T$ jest uniwersalna, a jej towarzysz
modelowy $T^*$ ma eliminację kwantyfikatorów,
to $T^*$ jest uzupełnieniem modelowym $T$.
\end{lemat}
\begin{proof}
Ustalmy $\mathfrak A\models T$.
W dowolnym modelu $T^*\cup\operatorname{Diag}(\mathfrak A)$
prawdziwość zdań bez kwantyfikatorów z nazwami elementów
$A$ jest wyznaczona przez diagram $\mathfrak A$:
dla każdego takiego zdania diagram zawiera je lub
jego negację.
Dowolne zdanie rozszerzonego języka zawiera tylko
skończenie wiele nowych nazw. Zastępujemy je
świeżymi zmiennymi wolnymi, stosujemy eliminację
kwantyfikatorów w $T^*$, a następnie podstawiamy
ponownie nazwy. Uzyskany równoważnik bez kwantyfikatorów
ma ustaloną wartość logiczną we wszystkich modelach
teorii z diagramem. To dowodzi jej zupełności.
\end{proof}

\begin{twierdzenie}[Eliminacja kwantyfikatorów nad niepustymi krotkami]
\label{thm:mod-mc-ap-qe-tuples}
Niech $T$ będzie teorią uniwersalną mającą towarzysza
modelowego $T^*$ i własność amalgamacji.
Dla każdej formuły $\varphi(\bar x)$ z ustaloną
niepustą listą zmiennych $\bar x$ zawierającą wszystkie
zmienne wolne $\varphi$ istnieje równoważnik bez
kwantyfikatorów $\psi(\bar x)$ modulo $T^*$.
\end{twierdzenie}
\begin{proof}
Niech $\mathfrak M,\mathfrak N\models T^*$ i niech
niepuste skończone krotki $\bar a,\bar b$ mają ten
sam typ bez kwantyfikatorów. Z lematu
\ref{lem:mod-qe-generated-substructures} wynika,
że generowane przez nie podstruktury są izomorficzne
przez odwzorowanie wysyłające $\bar a$ na $\bar b$.
Obie podstruktury są modelami $T$, ponieważ $T$
jest uniwersalna (twierdzenie
\ref{thm:mod-mc-universal-preservation}).
Traktujemy jedną z nich jako wspólną bazę amalgamacji,
z drugim zanurzeniem złożonym z izomorfizmu i inkluzji.
Własność amalgamacji daje wspólny model $T$,
w którym obrazy obu krotek są równe.
Rozszerzamy go do modelu $T^*$.
Zupełność modelowa $T^*$ zapewnia elementarność
obu zanurzeń modeli $T^*$ w ten model.
Stąd $\bar a,\bar b$ zgadzają się co do wszystkich
formuł przy ustalonej liście zmiennych.

Stosujemy argument zwartości z dowodu twierdzenia
\ref{thm:mod-qe-semantic} do ustalonej niepustej
listy $\bar x$. W dowodzie tym rozpatruje się jedynie
krotki tej długości, zatem zgodność, którą właśnie
wykazaliśmy, wystarcza do otrzymania równoważnika
bez kwantyfikatorów.
\end{proof}

\begin{twierdzenie}[Pełne kryterium amalgamacji]
\label{thm:mod-mc-ap-qe}
Niech $T$ będzie teorią uniwersalną mającą towarzysza
modelowego $T^*$. Następujące warunki są równoważne:
\begin{enumerate}[label=(\roman*)]
\item $T^*$ ma eliminację kwantyfikatorów
(z uwzględnieniem zdań);
\item $T$ ma własność amalgamacji i każde dwa modele
$T^*$ zgadzające się co do wszystkich zdań bez
kwantyfikatorów można zanurzyć w jeden model $T$.
\end{enumerate}
\end{twierdzenie}
\begin{proof}
Z (i) wynika własność amalgamacji na mocy lematu
\ref{lem:mod-mc-qe-completion} i twierdzenia
\ref{thm:mod-mc-completion-ap}.
Dwa modele $T^*$ zgodne co do zdań bez kwantyfikatorów
są wówczas elementarnie równoważne. Zatem wszystkie
zdania egzystencjalne prawdziwe w pierwszym modelu
są prawdziwe w drugim. Stosujemy lemat
\ref{lem:mod-mc-existential-type} do pustych krotek.
Otrzymujemy wspólny model $T^*$ (rozszerzenie
elementarne drugiego modelu), a więc także model $T$,
w którym zanurzają się oba modele.

Załóżmy (ii). Dla niepustych krotek zgodność
typów bez kwantyfikatorów daje zgodność wszystkich
formuł, jak w dowodzie twierdzenia
\ref{thm:mod-mc-ap-qe-tuples}.
Dla pustych krotek rozpatrujemy modele zgodne
co do zdań bez kwantyfikatorów. Z założenia zanurzamy
je w jeden model $T$, a ten rozszerzamy do modelu $T^*$.
Zupełność modelowa zapewnia elementarność
obu złożonych zanurzeń, więc oba modele zgadzają
się co do wszystkich zdań. Kryterium
\ref{thm:mod-qe-semantic} daje pełną eliminację.
\end{proof}

\begin{wniosek}[Dwa ważne przypadki]
\label{cor:mod-mc-ap-qe-standard}
Jeżeli $T$ jest uniwersalna i ma towarzysza modelowego
$T^*$, a ponadto język zawiera term domknięty lub
$T^*$ jest zupełna, to następujące warunki są
równoważne:
\begin{enumerate}[label=(\roman*)]
\item $T^*$ jest uzupełnieniem modelowym $T$;
\item $T$ ma własność amalgamacji;
\item $T^*$ ma eliminację kwantyfikatorów.
\end{enumerate}
\end{wniosek}
\begin{proof}
Z twierdzenia \ref{thm:mod-mc-completion-ap}
wynika równoważność (i) i (ii), a lemat
\ref{lem:mod-mc-qe-completion} daje (iii)$\Rightarrow$(i).
Pozostaje (ii)$\Rightarrow$(iii).
Jeżeli język zawiera term domknięty, modele
zgodne co do zdań bez kwantyfikatorów mają
izomorficzne niepuste podstruktury generowane
przez pustą krotkę (por. lemat
\ref{lem:mod-qe-generated-substructures}).
Są to modele $T$, więc można zastosować amalgamację.
Jeżeli natomiast $T^*$ jest zupełna, dowolne dwa
jej modele są elementarnie równoważne i można
je zanurzyć w jeden model $T^*$ przez lemat
\ref{lem:mod-mc-existential-type} dla pustych krotek.
W obu przypadkach spełniona jest dodatkowa hipoteza
twierdzenia \ref{thm:mod-mc-ap-qe}.
\end{proof}

\begin{przyklad}[Dlaczego trzeba uwzględnić zdania?]
\label{ex:mod-mc-ap-empty}
Rozważmy język z jednym jednoargumentowym symbolem
relacyjnym $P$ i teorię uniwersalną $T$ z aksjomatem
\[
\forall x\,\forall y\,(P(x)\leftrightarrow P(y)).
\]
Jej modele są niepustymi zbiorami, w których $P$
jest całym uniwersum lub zbiorem pustym.
Niech $T^*$ powstanie przez dodanie zdań
wymagających co najmniej $n$ różnych elementów
dla każdego $n\geq1$.

Każdy model $T$ można rozszerzyć do modelu $T^*$
przez dodanie elementów o tej samej wartości $P$.
Każdy model $T^*$ jest modelem $T$, więc obie teorie
mają te same konsekwencje uniwersalne na mocy lematu
\ref{lem:mod-mc-mutual-embedding}.
Każde zanurzenie modeli $T^*$ zachowuje wartość $P$:
uniwersum dziedziny jest niepuste, a $P$ jest
jednorodne w obu modelach. Po usunięciu symbolu $P$
jest to zanurzenie nieskończonych struktur języka
samej równości, a więc jest elementarne na mocy
twierdzenia \ref{thm:mod-qe-equality}.
Dla ustalonej wspólnej wartości $P$ każdą formułę
języka z $P$ można tłumaczyć na język samej równości,
zastępując atomy $P(x)$ odpowiednio przez $x=x$
lub $\neg(x=x)$. Indukcja po formułach daje zgodność
prawdziwości. Zatem zanurzenie jest elementarne
również w języku z $P$, co dowodzi zupełności
modelowej $T^*$.

Teoria $T$ ma nawet silną własność amalgamacji:
dwa rozszerzenia niepustego wspólnego modelu
muszą mieć tę samą wartość $P$.
Bierzemy rozłączne kopie poza wspólnym uniwersum,
łączymy je i nadajemy wszystkim elementom tę wartość.
Z twierdzenia \ref{thm:mod-mc-completion-ap}
$T^*$ jest uzupełnieniem modelowym $T$.
Nie ma jednak eliminacji kwantyfikatorów w sensie
przyjętym w tym skrypcie, obejmującym zdania.
W języku nie ma termów domkniętych ani relacji
o arności zero, więc każde zdanie bez kwantyfikatorów
jest kombinacją $\top,\bot$.
Zdanie $\exists x\,P(x)$ rozróżnia dwa modele $T^*$
o różnych wartościach $P$, a zatem nie ma takiego
równoważnika. Niepustość wspólnej bazy amalgamacji
jest tutaj istotna.
\end{przyklad}

\section{Przykłady amalgamacji i jej braku}
\begin{przyklad}[Grafy proste]
\label{ex:mod-mc-graphs-ap}
Teoria grafów prostych ma silną własność amalgamacji.
Dla dwóch grafów rozszerzających wspólny graf
$\mathfrak A$ zastępujemy je kopiami z $B\cap C=A$.
Na $D=B\cup C$ kładziemy $R^D=R^B\cup R^C$.
Relacja jest symetryczna i irrefleksyjna, ponieważ
każda z dwóch relacji ma te własności.
Jeżeli oba końce krawędzi z $R^C$ należą do $B$,
to należą do $A$, a zatem krawędź należy już
do $R^B$, ponieważ $\mathfrak A$ jest podstrukturą
obu grafów. W konsekwencji ograniczenie $R^D$
do $B^2$ to dokładnie $R^B$; analogicznie dla $C^2$.
Inkluzje są więc zanurzeniami i mają przecięcie $A$.
Towarzysz modelowy grafów prostych z przykładu
\ref{ex:mod-mc-companion-graphs} jest ich uzupełnieniem
modelowym na mocy twierdzenia
\ref{thm:mod-mc-completion-ap}.
\end{przyklad}

\begin{przyklad}[Porządki liniowe]
\label{ex:mod-mc-orders-ap}
Teoria porządków liniowych ma silną własność amalgamacji.
Niech $\mathfrak A$ będzie wspólną podstrukturą
porządków $\mathfrak B,\mathfrak C$, z $B\cap C=A$.
Na $D=B\cup C$ rozważmy relację będącą przechodnim
domknięciem sumy obu ścisłych porządków.
Każde porównanie w tym domknięciu odpowiada skończonej
niepustej drodze rosnących kroków, z których każdy
pochodzi z jednego z porządków.

Jeżeli końce takiej drogi należą do $B$, każdy
maksymalny odcinek drogi przechodzący przez
$C\setminus A$ ma końce w $A$: przejście między
$B$ i $C$ jest możliwe tylko we wspólnym uniwersum.
Odcinek ten zastępujemy pojedynczym krokiem w $B$,
ponieważ przechodniość porządku $C$ daje porównanie
jego końców, a oba porządki zgadzają się na $A$.
Otrzymana niepusta droga w $B$ daje porównanie
jej końców w pierwotnym porządku $B$.
Analogicznie dla końców należących do $C$.
W szczególności nie ma drogi od elementu do niego
samego, gdyż każdy element należy do $B$ lub $C$,
a pierwotne porządki są irrefleksyjne.
Przechodnie domknięcie jest zatem ścisłym porządkiem
częściowym, którego ograniczenia do $B,C$ są
pierwotnymi porządkami.

Każdy ścisły porządek częściowy można rozszerzyć
do ścisłego porządku liniowego na tym samym zbiorze.
Dla przypomnienia stosujemy lemat Zorna do zbioru
ścisłych porządków częściowych rozszerzających
ustalony porządek, uporządkowanego przez inkluzję.
Suma łańcucha jest ponownie ścisłym porządkiem
częściowym (każde dwa kroki leżą w jednym członie).
Gdyby maksymalny porządek nie porównywał $x,y$,
dodanie $x<y$ i wzięcie przechodniego domknięcia
nie tworzyłoby cyklu: z cyklu zawierającego nowy
krok można wyodrębnić drogę od $y$ do $x$ po starych
krokach, sprzeczną z nieporównywalnością.
Otrzymalibyśmy właściwe rozszerzenie porządku,
co przeczy maksymalności.

Tak uzyskany porządek liniowy na $D$ nadal ma
pierwotne ograniczenia do $B,C$: każde porównanie
różnych elementów w każdym z tych zbiorów było
już rozstrzygnięte. Inkluzje są zanurzeniami,
a ich obrazy przecinają się dokładnie w $A$.
Towarzysz teorii porządków liniowych, czyli DLO
(przykład \ref{ex:mod-mc-companion-dlo}), jest
zatem jej uzupełnieniem modelowym.
\end{przyklad}

\begin{przyklad}[Towarzysz bez uzupełnienia modelowego]
\label{ex:mod-mc-no-ap}
Niech $S=\operatorname{Th}(\R,0,1,+,-,\cdot)$
w języku pierścieni, a $T=S_{\forall}$.
Z przykładu \ref{ex:mod-mc-no-qe} wiemy, że $S$
jest modelowo zupełna, ale nie ma eliminacji
kwantyfikatorów. Z definicji $S$ jest towarzyszem
modelowym $T$, ponieważ $T_{\forall}=S_{\forall}$.
Ponieważ język zawiera stałą $0$, z wniosku
\ref{cor:mod-mc-ap-qe-standard} wynika, że $T$
nie ma własności amalgamacji i $S$ nie jest
uzupełnieniem modelowym $T$.

Można także wskazać konkretną przeszkodę amalgamacji.
Niech $\mathfrak A$ będzie pierścieniem
$\mathbb Z[X]$. Wybieramy dodatnią liczbę przestępną
$t\in\R$ jak w przykładzie \ref{ex:mod-mc-no-qe}.
Podstawienia $X\mapsto t$ i $X\mapsto -t$ są
zanurzeniami $\mathfrak A$ w dwa egzemplarze $\R$,
ponieważ żaden niezerowy wielomian całkowity
nie zeruje się w $t$ ani w $-t$.
Z twierdzenia \ref{thm:mod-mc-universal-embedding}
$\mathfrak A\models T$, a oba egzemplarze $\R$
są modelami $T$. Gdyby istniała amalgamacja w
modelu $\mathfrak D\models T$, można by rozszerzyć
$\mathfrak D$ do modelu $S$.
W tym modelu wspólny obraz $X$ byłby niezerowym
kwadratem, ponieważ w pierwszym egzemplarzu $\R$
jest nim $t$. Jego przeciwny element byłby również
niezerowym kwadratem, ponieważ w drugim egzemplarzu
obrazem $-X$ jest $t$. Jest to niemożliwe w modelu
$S$, który, jak wykazano w przykładzie
\ref{ex:mod-mc-no-qe}, ma porządek ciała rzeczywiście
domkniętego. W takim porządku niezerowy kwadrat
jest dodatni, a jego przeciwny element ujemny.
\end{przyklad}

Zupełność modelowa, towarzysze i uzupełnienia modelowe
pozwalają zatem oddzielić zachowanie formuł w modelach
danej teorii od możliwości rozszerzania jej modeli
oraz uzgadniania rozszerzeń nad wspólną bazą.
W dalszych rozdziałach wykorzystamy te pojęcia
do badania modeli pierwszych, kategoryczności
i własności zbiorów definiowalnych.



